%========================================================
% 자율학습 관리 시스템 사용설명서
% 컴파일: xelatex manual.tex (2회 이상 실행 권장)
%========================================================
\documentclass[a4paper,12pt]{article}

%--- 기본 패키지 ---
\usepackage{kotex}
\usepackage[a4paper, top=2.5cm, bottom=2.5cm, left=2.8cm, right=2.5cm]{geometry}
\usepackage{xcolor}
\usepackage{hyperref}
\usepackage{booktabs}
\usepackage{tabularx}
\usepackage{array}
\usepackage{enumitem}
\usepackage{fancyhdr}
\usepackage{titlesec}
\usepackage{tcolorbox}
\usepackage{listings}
\usepackage{graphicx}
\usepackage{float}
\usepackage{mdframed}
\usepackage{fontawesome5}

%--- 색상 정의 ---
\definecolor{primary}{HTML}{0d6efd}
\definecolor{success}{HTML}{198754}
\definecolor{warning}{HTML}{ffc107}
\definecolor{danger}{HTML}{dc3545}
\definecolor{orange}{HTML}{fd7e14}
\definecolor{info}{HTML}{0dcaf0}
\definecolor{secondary}{HTML}{6c757d}
\definecolor{dark}{HTML}{343a40}
\definecolor{codebg}{HTML}{f8f9fa}
\definecolor{codefg}{HTML}{212529}
\definecolor{notecolor}{HTML}{cfe2ff}
\definecolor{warncolor}{HTML}{fff3cd}
\definecolor{tipcolor}{HTML}{d1e7dd}
\definecolor{purple}{HTML}{6f42c1}

%--- 하이퍼링크 설정 ---
\hypersetup{
    colorlinks=true,
    linkcolor=primary,
    urlcolor=primary,
    pdftitle={자율학습 관리 시스템 사용설명서},
    pdfauthor={시스템 관리자},
}

%--- 헤더·푸터 ---
\pagestyle{fancy}
\fancyhf{}
\fancyhead[L]{\small 자율학습 관리 시스템}
\fancyhead[R]{\small \leftmark}
\fancyfoot[C]{\thepage}
\renewcommand{\headrulewidth}{0.4pt}

%--- 섹션 스타일 ---
\titleformat{\section}
  {\Large\bfseries\color{primary}}
  {\thesection.}{0.5em}{}[\titlerule]

\titleformat{\subsection}
  {\large\bfseries\color{primary!80!black}}
  {\thesubsection}{0.5em}{}

\titleformat{\subsubsection}
  {\normalsize\bfseries}
  {\thesubsubsection}{0.5em}{}

%--- 코드 스타일 ---
\lstset{
    backgroundcolor=\color{codebg},
    basicstyle=\ttfamily\small,
    keywordstyle=\color{primary}\bfseries,
    commentstyle=\color{gray},
    stringstyle=\color{success},
    frame=single,
    framerule=0.4pt,
    rulecolor=\color{gray!40},
    breaklines=true,
    breakatwhitespace=true,
    showstringspaces=false,
    columns=flexible,
    keepspaces=true,
    xleftmargin=6pt,
    xrightmargin=6pt,
    aboveskip=8pt,
    belowskip=8pt,
}

%--- 알림 박스 ---
\tcbuselibrary{skins, breakable}

\newtcolorbox{notebox}{
    enhanced, breakable,
    colback=notecolor, colframe=primary,
    leftrule=4pt, rightrule=0pt, toprule=0pt, bottomrule=0pt,
    arc=0pt, outer arc=0pt,
    title={\faInfoCircle\; 참고},
    fonttitle=\bfseries\small,
    coltitle=primary,
    attach boxed title to top left={yshift=-2mm, xshift=4mm},
    boxed title style={colback=white, colframe=primary},
}

\newtcolorbox{warnbox}{
    enhanced, breakable,
    colback=warncolor, colframe=warning!80!black,
    leftrule=4pt, rightrule=0pt, toprule=0pt, bottomrule=0pt,
    arc=0pt, outer arc=0pt,
    title={\faExclamationTriangle\; 주의},
    fonttitle=\bfseries\small,
    coltitle=warning!60!black,
    attach boxed title to top left={yshift=-2mm, xshift=4mm},
    boxed title style={colback=white, colframe=warning!80!black},
}

\newtcolorbox{tipbox}{
    enhanced, breakable,
    colback=tipcolor, colframe=success,
    leftrule=4pt, rightrule=0pt, toprule=0pt, bottomrule=0pt,
    arc=0pt, outer arc=0pt,
    title={\faCheckCircle\; 팁},
    fonttitle=\bfseries\small,
    coltitle=success,
    attach boxed title to top left={yshift=-2mm, xshift=4mm},
    boxed title style={colback=white, colframe=success},
}

%--- 헤더 높이 조정 ---
\setlength{\headheight}{14pt}

%--- 메뉴 경로 표시 ---
\newcommand{\menu}[1]{\textcolor{primary}{\textbf{#1}}}
\newcommand{\myurl}[1]{\texttt{\textcolor{primary}{#1}}}
\newcommand{\key}[1]{\fbox{\small\texttt{#1}}}
% badge: 기본 흰 글자 / badgedark: 어두운 글자(밝은 배경용 - warning, info 등)
\newcommand{\badge}[2]{\colorbox{#1}{\small\textcolor{white}{\textbf{#2}}}}
\newcommand{\badgedark}[2]{\colorbox{#1}{\small\textcolor{codefg}{\textbf{#2}}}}

%========================================================
\begin{document}
%========================================================

%--- 표지 ---
\begin{titlepage}
\centering
\vspace*{3cm}

{\Huge\bfseries\color{primary} 자율학습 관리 시스템}\\[0.8em]
{\LARGE 사용설명서}\\[3cm]

\begin{tcolorbox}[
    enhanced, width=0.8\textwidth,
    colback=primary!10, colframe=primary,
    arc=8pt,
]
\centering
\Huge
고등학생 자율학습 신청·출석·기록 통합 관리 플랫폼\\[0.5em]
\large
Flask 기반 웹 애플리케이션 \\
QR코드 입실·퇴실·조퇴처리·월별 신청·교사 승인 시스템
\end{tcolorbox}

\vfill

{\large \today}\\[0.5em]
{\normalsize 본 문서는 시스템 관리자 및 사용자를 위한 공식 사용설명서입니다.}
\end{titlepage}

%--- 목차 ---
\tableofcontents
\newpage

%========================================================
\section{시스템 개요}
%========================================================

자율학습 관리 시스템은 고등학교에서 방과 후 자율학습을 체계적으로 운영하기 위한 웹 기반 통합 관리 플랫폼입니다.

\subsection{주요 기능}

\begin{itemize}[leftmargin=1.5em]
    \item \textbf{월별 자습 신청}: 학생이 달력 형태의 UI에서 원하는 날짜·교시를 선택하여 신청
    \item \textbf{QR코드 입실}: 자습실 입구의 입실 QR코드를 스캔하여 자동 출석 처리
    \item \textbf{QR코드 퇴실·조퇴 처리}: 자습 종료 전 퇴실 시 퇴실 QR코드를 스캔하면 자동 조퇴 처리; 교사가 사유 확인 후 출석인정으로 변환 가능
    \item \textbf{출석 관리}: 교사가 날짜·학년·공간별로 출석 현황을 조회·수정
    \item \textbf{학습 기록}: 학생이 과목별 학습 시간과 메모를 직접 기록
    \item \textbf{참여 통계}: 월별 신청 대비 출석률 상위 학생 조회
    \item \textbf{자습 공간 관리}: 독서실·열람실 등 자습 공간 등록 및 학생 배정
    \item \textbf{교사 승인 시스템}: 관리자가 교사 계정을 승인한 후에만 시스템 접근 허용
    \item \textbf{좌석 배치도 시각화}: 자습실별 학생 배치를 드래그 앤 드롭으로 조정·공유, 날짜·교시별 실시간 출결 현황 표시
    \item \textbf{수동 자리 배정}: 교사가 각 번호 슬롯에 학생을 드롭다운으로 직접 지정 (임의 배정 대안)
    \item \textbf{자동 지각/결석 처리}: 신청 학생 중 미처리 학생을 시간 기준으로 일괄 처리
    \item \textbf{출결 수정 이력}: 출결 변경 내역(수동·자동)을 날짜별로 조회
    \item \textbf{학생 개인 리포트}: 월별 신청·출석·참여율·학습시간 요약
    \item \textbf{Excel 내보내기}: 출석 현황 및 참여 통계를 엑셀 파일로 다운로드
    \item \textbf{학생 출결 자기확인}: 학생이 월별 본인 출결 이력(출석·지각·결석·참여율)을 직접 조회
    \item \textbf{CSRF 보안 보호}: 모든 폼 및 API 요청에 위변조 방지 토큰 적용
    \item \textbf{새 학년도 초기화}: Excel 백업 후 학생 데이터 일괄 삭제, 교사·설정 보존
\end{itemize}

\subsection{사용자 역할}

\begin{center}
\begin{tabularx}{\textwidth}{lX}
\toprule
\textbf{역할} & \textbf{권한} \\
\midrule
\badge{danger}{관리자} & 교사 계정 승인·거부·권한취소, 학생·교사 계정 삭제, 비밀번호 초기화, 관리자 비밀번호 변경 \\[0.3em]
\badge{success}{교사} & 자습 시간 설정, 공휴일 등록, 자습 공간 관리, 학생 배정, 출석 관리, 통계 조회 \\[0.3em]
\badge{primary}{학생} & 자습 신청, QR 출석, 학습 기록 \\
\bottomrule
\end{tabularx}
\end{center}

\vspace{1em}
\begin{notebox}
  \textbf{교사 계정}은 회원가입 후 관리자의 승인이 있어야만 시스템에 로그인할 수 있습니다.
  \textbf{학생 계정}은 회원가입 즉시 로그인이 가능합니다.
\end{notebox}

%========================================================
\section{시스템 요구사항 및 설치}
%========================================================

\subsection{소프트웨어 요구사항}

\begin{center}
\begin{tabular}{ll}
\toprule
\textbf{항목} & \textbf{버전} \\
\midrule
Python & 3.10 이상 \\
운영체제 & Windows 10/11, macOS, Linux \\
웹 브라우저 & Chrome, Edge, Firefox (최신 버전) \\
\bottomrule
\end{tabular}
\end{center}

\subsection{파이썬 패키지 설치}

터미널(명령 프롬프트)을 열고 프로젝트 폴더로 이동한 후 아래 명령을 실행합니다.

\begin{lstlisting}[language=bash, caption=필수 패키지 설치]
pip install flask flask-sqlalchemy flask-login werkzeug waitress openpyxl flask-wtf
\end{lstlisting}

\subsection{프로젝트 파일 구조}

\begin{lstlisting}[caption=파일 구조]
self_study/
├── app.py                  # Flask 앱 진입점
├── auth.py                 # 로그인·회원가입 처리
├── models.py               # 데이터베이스 모델
├── routes_student.py       # 학생 기능 라우트
├── routes_teacher.py       # 교사 기능 라우트
├── routes_admin.py         # 관리자 기능 라우트
├── reset_db.py             # 데이터베이스 초기화 스크립트
├── instance/
│   ├── self_study.db       # SQLite 데이터베이스
│   └── secret_key.txt      # 세션 암호화 키 (자동 생성)
└── templates/
    ├── base.html           # 공통 레이아웃
    ├── login.html          # 로그인 페이지
    ├── register.html       # 회원가입 페이지
    ├── pending_approval.html  # 교사 승인 대기 화면
    ├── admin/
    │   ├── dashboard.html  # 관리자 대시보드
    │   ├── teachers.html   # 교사 승인 관리
    │   ├── users.html      # 사용자 관리 (삭제·비밀번호 초기화)
    │   └── change_password.html  # 비밀번호 변경
    ├── student/
    │   ├── dashboard.html  # 학생 대시보드
    │   ├── apply.html      # 자습 신청 달력
    │   └── study_log.html  # 학습 기록
    └── teacher/
        ├── dashboard.html     # 교사 대시보드
        ├── students.html      # 학생 목록
        ├── applications.html  # 자습 신청 현황
        ├── attendance.html    # 출석 관리
        ├── statistics.html    # 참여 통계
        ├── settings.html      # 시스템 설정
        ├── qr_display.html    # QR코드 표시
        ├── seat_layout.html   # 좌석 배치도 (드래그 앤 드롭)
        ├── manual_seats.html  # 수동 자리 배정
        └── student_report.html # 학생 개인 리포트
\end{lstlisting}

%========================================================
\section{시스템 시작 (최초 설정)}
%========================================================

\subsection{데이터베이스 초기화}

시스템을 처음 사용하거나 초기화가 필요할 때는 다음 명령을 실행합니다.

\begin{lstlisting}[language=bash, caption=데이터베이스 초기화]
cd self_study
python reset_db.py
\end{lstlisting}

\begin{warnbox}
  \texttt{reset\_db.py}를 실행하면 \textbf{기존의 모든 데이터(학생, 교사, 출석, 신청 내역)}가
  삭제됩니다. 운영 중에는 절대 실행하지 마십시오.
\end{warnbox}

\subsection{서버 실행}

\begin{lstlisting}[language=bash, caption=서버 실행]
python app.py
\end{lstlisting}

서버가 정상적으로 시작되면 터미널에 다음과 같이 출력됩니다.

\begin{lstlisting}[caption=정상 실행 출력 예시]
==================================================
  자율학습 관리 시스템
  http://192.168.x.x:5000  <- 학생/교사 모두 이 주소로 접속
==================================================
\end{lstlisting}

\begin{notebox}
  출력되는 IP 주소가 서버 PC의 실제 네트워크 주소입니다.
  교사와 학생 모두 이 주소를 브라우저에 입력하여 접속합니다.
\end{notebox}

\subsection{관리자 계정 자동 생성}

\textbf{최초 실행 시}, 데이터베이스에 관리자 계정이 없으면 기본 관리자 계정이 자동으로 생성되며
터미널에 아래와 같이 출력됩니다.

\begin{lstlisting}[caption=관리자 계정 자동 생성 메시지]
==================================================
  [관리자 계정 자동 생성]
  아이디: admin
  비밀번호: Admin1234!
  ※ 첫 로그인 후 반드시 비밀번호를 변경하세요!
==================================================
\end{lstlisting}

\begin{warnbox}
  이 메시지는 관리자 계정이 없을 때만 출력됩니다.
  비밀번호를 메모한 후, 첫 로그인 즉시 \textbf{비밀번호 변경}을 진행하십시오.
\end{warnbox}

\subsection{브라우저 접속}

서버 실행 후 터미널에 출력된 IP 주소를 웹 브라우저에 입력하여 접속합니다.

\begin{center}
\large\texttt{\textcolor{primary}{http://[서버 PC의 IP 주소]:5000}}
\end{center}

서버 PC의 IP 주소는 다음 명령으로 확인합니다.

\begin{lstlisting}[language=bash, caption=서버 PC IP 주소 확인]
ipconfig
# "이더넷 어댑터" 또는 "무선 LAN 어댑터"의 IPv4 주소 확인
\end{lstlisting}

\begin{tipbox}
  서버는 \texttt{0.0.0.0}으로 바인딩되어 있어 교사와 학생 모두
  같은 IP 주소로 접속합니다. 서버 PC 자체에서 접속할 때도
  \texttt{127.0.0.1:5000} 대신 실제 IP 주소를 사용하는 것을 권장합니다.
\end{tipbox}

\subsection{서버 자동 시작 설정 (Windows)}

서버 PC가 재부팅될 때 \texttt{app.py}가 자동으로 실행되도록 설정할 수 있습니다.
본 시스템은 Flask 내장 서버 대신 \textbf{Waitress} WSGI 서버를 사용합니다.
Waitress는 Wi-Fi 끊김·재연결 후에도 소켓을 자동 복구하여 안정적으로 동작합니다.

\subsubsection{사전 준비}

\begin{enumerate}[leftmargin=1.5em]
    \item \textbf{Waitress 설치} --- 아직 설치하지 않았다면 실행합니다.
\begin{lstlisting}[language=bash]
pip install waitress
\end{lstlisting}

    \item \textbf{logs 폴더 생성} --- 서버 로그를 저장할 폴더를 만듭니다.
\begin{lstlisting}[language=bash]
mkdir c:\python_projects\self_study\logs
\end{lstlisting}
\end{enumerate}

\subsubsection{배치 파일(bat)}

프로젝트 폴더에 \texttt{start\_self\_study.bat} 파일이 생성되어 있습니다.
서버 실행 로그는 \texttt{logs\textbackslash app.log}에 저장됩니다.

\subsubsection{작업 스케줄러 등록}

\textbf{방법 A --- 명령어로 등록 (권장)}

관리자 권한 명령 프롬프트(cmd)에서 다음 명령을 실행합니다.

\begin{lstlisting}[language=bash, caption=작업 스케줄러 명령어 등록]
schtasks /create /tn "SelfStudyApp" ^
  /tr "c:\python_projects\self_study\start_self_study.bat" ^
  /sc onstart /ru SYSTEM /rl HIGHEST /f
\end{lstlisting}

\textbf{방법 B --- GUI로 등록}

\begin{enumerate}[leftmargin=1.5em]
    \item \key{Win}+\key{R} → \texttt{taskschd.msc} 입력 후 실행
    \item 오른쪽 패널에서 \menu{작업 만들기} 클릭
    \item \textbf{일반} 탭: 이름 \texttt{SelfStudyApp} 입력,
          ``가장 높은 수준의 권한으로 실행'' 체크,
          실행 조건을 \menu{사용자가 로그온되어 있는지 여부와 관계없이 실행}으로 변경
    \item \textbf{트리거} 탭 → \menu{새로 만들기} → \menu{시작 시} 선택 → 확인
    \item \textbf{동작} 탭 → \menu{새로 만들기} → 프로그램/스크립트:\\
          \texttt{c:\textbackslash python\_projects\textbackslash self\_study\textbackslash start\_self\_study.bat}
    \item \menu{확인} → 관리자 비밀번호 입력
\end{enumerate}

\subsubsection{실패 시 자동 재시작 설정}

작업 스케줄러에서 \texttt{SelfStudyApp} 작업 속성 열기 →
\textbf{설정} 탭에서 아래와 같이 설정합니다.

\begin{center}
\begin{tabular}{ll}
\toprule
\textbf{항목} & \textbf{권장 값} \\
\midrule
작업 실패 시 다시 시작 간격 & 1분 \\
다시 시작 시도 횟수 & 10회 \\
작업이 계속 실행되면 중지 시간 & 사용 안 함 \\
\bottomrule
\end{tabular}
\end{center}

\subsubsection{등록 확인}

재부팅 후 서버가 실행 중인지 확인합니다.

\begin{lstlisting}[language=bash, caption=실행 확인]
tasklist | findstr python
\end{lstlisting}

또는 브라우저에서 \texttt{http://[서버IP]:5000} 접속이 되면 정상입니다.

\begin{tipbox}
  \texttt{logs\textbackslash app.log}: 서버 표준 출력(시작 메시지·요청·예외).\\
  \texttt{logs\textbackslash audit.log}: 인증·관리자 작업 등 추적 대상 이벤트
  (5MB 단위 회전, 최대 5개 보관).
\end{tipbox}

\subsection{감사 로그 (logs/audit.log)}

운영 책임 추적을 위해 다음 이벤트가 \texttt{logs\textbackslash audit.log}에
구조화된 형태로 기록됩니다.

\begin{center}
\begin{tabularx}{\textwidth}{lX}
\toprule
\textbf{이벤트} & \textbf{기록 시점} \\
\midrule
\texttt{auth.login\_success}        & 로그인 성공 \\
\texttt{auth.login\_fail}           & 로그인 실패 (WARNING 레벨) \\
\texttt{auth.login\_denied}         & 미승인 교사 로그인 시도 \\
\texttt{auth.logout}                & 로그아웃 \\
\texttt{admin.pw\_reset}            & 관리자가 학생·교사 비밀번호 초기화 (WARNING) \\
\texttt{admin.account\_delete}      & 관리자가 계정 삭제 (WARNING) \\
\texttt{admin.db\_restore}          & 관리자가 DB 파일 복원 (WARNING) \\
\texttt{admin.excel\_restore}       & 관리자가 Excel 백업 복원 (WARNING) \\
\texttt{admin.new\_year\_reset}     & 새 학년도 초기화 (WARNING) \\
\texttt{system.admin\_account\_created} & 최초 부팅 시 admin 계정 자동 생성 \\
\texttt{system.scheduler\_started}  & 야간 자동 처리 스케줄러 시작 \\
\texttt{system.auto\_early\_leave}  & 자정 자동 조퇴 처리 (퇴실 미확인) \\
\texttt{system.constraints\_missing} & 부팅 시 DB 무결성 제약 누락 감지 (WARNING). 마이그레이션 미실행 의미 \\
\texttt{system.session\_token\_backfilled} & 부팅 시 NULL session\_token 발견 → UUID 자동 채움 (WARNING) \\
\texttt{system.session\_token\_check\_failed} & session\_token 컬럼 자체가 없음 → \texttt{migrate.py} 실행 필요 (ERROR) \\
\bottomrule
\end{tabularx}
\end{center}

기록 형식:

\begin{lstlisting}
2026-04-23 21:00:01 [INFO] self_study.audit: auth.login_success username=teacher01 role=teacher ip=192.168.0.42
2026-04-23 21:05:18 [WARNING] self_study.audit: admin.pw_reset admin=admin target=10301 target_role=student
\end{lstlisting}

\textbf{보안}: 평문 비밀번호는 절대 기록되지 않습니다.
PW 초기화 이벤트는 ``누가 누구를 초기화했다''만 남깁니다.

%========================================================
\section{관리자 기능}
%========================================================

관리자는 시스템 전체를 관장하며, 교사 계정의 승인 권한과
학생·교사 계정의 삭제 및 비밀번호 초기화 권한을 가집니다.

\subsection{관리자 로그인}

\begin{enumerate}[leftmargin=1.5em]
    \item 브라우저에서 \texttt{http://[서버IP]:5000}에 접속합니다.
    \item \textbf{아이디}: \texttt{admin}
    \item \textbf{비밀번호}: \texttt{Admin1234!} (최초 기본값)
    \item \menu{로그인} 버튼을 클릭합니다.
\end{enumerate}

로그인 성공 시 \textbf{관리자 대시보드}로 이동합니다.

\subsection{관리자 대시보드}

대시보드에서는 다음 정보를 한눈에 확인할 수 있습니다.

\begin{itemize}[leftmargin=1.5em]
    \item \textbf{전체 학생 수}: 등록된 학생 계정 수
    \item \textbf{승인된 교사 수}: 접근 권한이 있는 교사 수
    \item \textbf{승인 대기 교사 수}: 아직 승인되지 않은 교사 계정 수
\end{itemize}

대시보드 하단의 버튼으로 바로 이동할 수 있습니다.

\begin{center}
\begin{tabular}{ll}
\toprule
\textbf{버튼} & \textbf{이동 대상} \\
\midrule
\menu{사용자 관리} & 학생·교사 계정 목록, 삭제, 비밀번호 초기화 \\
\menu{교사 승인 관리} & 교사 계정 승인·거부·권한 취소 \\
\menu{시스템 설정} & 사전 입실·퇴실 grace·지각 임계·신청 마감일·비밀번호 정책 등 운영 정책값 일괄 관리 \\
\menu{비밀번호 변경} & 관리자 비밀번호 변경 \\
\bottomrule
\end{tabular}
\end{center}

\begin{tipbox}
  승인 대기 중인 교사가 있을 경우, 경고 알림이 표시되며 \menu{바로 승인하기} 버튼이 나타납니다.
\end{tipbox}

\subsection{사용자 관리 --- 학생·교사 계정 삭제 및 비밀번호 초기화}

상단 네비게이션에서 \menu{사용자 관리}를 클릭하면 전체 학생·교사 목록이 표시됩니다.

\subsubsection{목록 조회 및 검색}

\begin{itemize}[leftmargin=1.5em]
    \item \textbf{역할 필터 탭}: \menu{전체}, \menu{학생}, \menu{교사} 탭으로 전환합니다.
    \item \textbf{검색}: 이름 또는 아이디의 일부를 입력하고 \menu{검색} 버튼을 클릭합니다.
    \item 검색 결과를 초기화하려면 \menu{초기화} 버튼을 클릭합니다.
\end{itemize}

표에는 ID, 역할, 이름, 아이디, 학년·반(학생만), 승인 상태(교사만)가 표시됩니다.

\subsubsection{비밀번호 초기화}

\begin{enumerate}[leftmargin=1.5em]
    \item 초기화할 사용자의 행에서 \menu{PW 초기화} 버튼을 클릭합니다.
    \item 확인 창에서 \key{확인}을 선택합니다.
    \item 화면 상단에 \textbf{임시 비밀번호}가 표시됩니다.
    \item 임시 비밀번호를 해당 학생·교사에게 전달하고, 로그인 후 즉시 변경하도록 안내합니다.
\end{enumerate}

\begin{warnbox}
  임시 비밀번호는 \textbf{이 화면에서 한 번만 표시}됩니다. 반드시 메모하거나
  즉시 해당 사용자에게 전달하십시오. 페이지를 이동하면 다시 확인할 수 없습니다.
\end{warnbox}

\begin{notebox}
  임시 비밀번호는 영문+숫자 조합 8자리로 자동 생성됩니다.
  비밀번호 규칙(8자 이상, 영문+숫자 포함)을 만족하므로 학생·교사가 즉시 로그인할 수 있습니다.
\end{notebox}

\subsubsection{계정 삭제}

\begin{enumerate}[leftmargin=1.5em]
    \item 삭제할 사용자의 행에서 \menu{삭제} 버튼을 클릭합니다.
    \item 확인 창에 삭제 범위(출석, 신청, 학습기록 등 모든 데이터)가 표시됩니다.
    \item \key{확인}을 선택하면 즉시 삭제됩니다.
\end{enumerate}

\begin{warnbox}
  계정을 삭제하면 해당 사용자의 \textbf{자습 신청, 출석 기록, 학습 기록, 공간 배정}이
  모두 함께 삭제됩니다. 삭제 후에는 복구할 수 없으므로 신중하게 사용하십시오.
\end{warnbox}

\begin{notebox}
  \texttt{admin} 계정(관리자)은 사용자 관리 목록에서 표시되지 않으며, 삭제·초기화가 불가능합니다.
\end{notebox}

\subsection{교사 계정 승인}

\begin{enumerate}[leftmargin=1.5em]
    \item 상단 네비게이션에서 \menu{교사 승인}을 클릭합니다.
    \item \menu{승인 대기} 탭을 선택하면 승인을 기다리는 교사 목록이 표시됩니다.
    \item 해당 교사 행의 \menu{승인} 버튼을 클릭합니다.
    \item 확인 창에서 \key{확인}을 누르면 교사 계정이 활성화됩니다.
\end{enumerate}

\begin{center}
\begin{tabular}{lll}
\toprule
\textbf{버튼} & \textbf{상태} & \textbf{동작} \\
\midrule
\badge{success}{승인} & 대기 → 활성 & 교사가 로그인하여 시스템 사용 가능 \\
\badge{danger}{거부} & 대기 → 삭제 & 계정 자체를 데이터베이스에서 삭제 \\
\badge{danger}{권한 취소} & 활성 → 대기 & 이미 승인된 교사의 접근 권한을 중지 \\
\bottomrule
\end{tabular}
\end{center}

\subsection{새 학년도 초기화}

매년 학년도 종료(보통 2월 말) 후 학생 데이터를 초기화하여 새 학년도를 시작할 수 있습니다.

\begin{warnbox}
  초기화를 실행하면 \textbf{학생 데이터가 영구적으로 삭제}되며 복구할 수 없습니다.
  반드시 \textbf{백업을 먼저 다운로드}한 후 진행하십시오.
\end{warnbox}

\subsubsection{삭제 및 보존 범위}

\begin{center}
\begin{tabular}{ll}
\toprule
\textbf{삭제되는 데이터} & \textbf{보존되는 데이터} \\
\midrule
학생 계정 전체 & 관리자·교사 계정 \\
자습 신청 기록 전체 & 자습실 설정 \\
출결 기록 전체 & 자습 시간 설정 \\
학습 기록 전체 & 공휴일 설정 \\
출결 수정 이력 전체 & \\
\bottomrule
\end{tabular}
\end{center}

\subsubsection{초기화 절차 (2단계)}

\begin{enumerate}[leftmargin=1.5em]
    \item 관리자 대시보드에서 \menu{새 학년도 초기화} 버튼을 클릭합니다.
    \item 현황 카드에서 삭제될 데이터 건수를 확인합니다.
    \item \textbf{1단계}: \menu{📥 백업 Excel 다운로드} 버튼을 클릭하여 전체 데이터를 백업합니다.
    \item \textbf{2단계}: 확인 입력란에 \texttt{새학년도초기화}를 정확히 입력합니다.
    \begin{itemize}
        \item 문구가 일치하면 \menu{새 학년도 초기화 실행} 버튼이 활성화됩니다.
    \end{itemize}
    \item 버튼 클릭 후 최종 확인 창에서 \key{확인}을 선택합니다.
    \item 초기화 완료 후 대시보드로 이동하며 결과가 표시됩니다.
\end{enumerate}

\subsubsection{백업 Excel 구성}

\menu{📥 백업 Excel 다운로드}를 클릭하면 다음 11개 시트로 구성된 Excel 파일이 생성됩니다.

\begin{center}
\begin{tabularx}{\textwidth}{llX}
\toprule
\textbf{시트} & \textbf{복원} & \textbf{내용} \\
\midrule
학생명단     & ✓      & 학번, 이름, 학년, 반, 성별, 아이디 \\
출결요약     & 참고용  & 학생별 월별 출석/지각/결석/신청 횟수·참여율 \\
학습기록     & ✓      & 학생별 과목·학습시간·메모 전체 \\
자습신청     & ✓      & 학생별 날짜·교시 신청 목록 전체 \\
출결상세     & ✓      & 날짜·교시·상태·출석·퇴실시각·조퇴사유 원시 데이터 \\
교사명단     & ✓      & 아이디, 이름, 승인여부, 담당학년 \\
관리자       & ✓      & 아이디, 이름 \\
자습시간설정 & ✓      & 요일별(월\textasciitilde금·토·공휴일) 교시 시작·종료 시각, 활성화 여부 \\
공휴일       & ✓      & 날짜, 공휴일명 \\
자습실목록   & ✓      & 자습실명, 전체·남·여 정원, 활성화 여부 \\
자습실배정   & ✓      & 학생별 자습실명, 좌석번호, 배치도 위치(X/Y) \\
\bottomrule
\end{tabularx}
\end{center}

파일명은 \texttt{자율학습\_백업\_YYYYMMDD\_HHMM.xlsx} 형식으로 저장됩니다.

\begin{notebox}
  보안상 Excel 백업에는 \textbf{비밀번호 해시·QR 토큰이 포함되지 않습니다}.
  비밀번호를 그대로 보존하면서 복원하려면 \menu{⬇ DB 파일 백업 (.db)}을 사용하십시오.
  Excel 복원 경로는 \textbf{신규 학년도 데이터 이동용}으로, 복원 시 모든 신규 계정에
  랜덤 임시 비밀번호가 발급되며 그 명단이 별도 Excel로 다운로드됩니다.
\end{notebox}

\subsubsection{초기화 후 작업}

\begin{enumerate}[leftmargin=1.5em]
    \item 새 학년도 학생들이 회원가입합니다.
    \item 교사가 자습실에 학생을 배정합니다.
    \item 학생들이 월별 자습 신청을 시작합니다.
\end{enumerate}

\begin{notebox}
  초기화는 언제든지 실행할 수 있으나, 학년도 중간에 실행하면
  해당 연도 데이터가 모두 사라집니다. 반드시 학년도 종료 직후에 실행하십시오.
\end{notebox}

\subsection{백업 파일로 복원}

새 학년도 초기화 후 또는 DB 손상 시 백업 Excel 파일로 데이터를 복원합니다.

\subsubsection{복원 절차}

\begin{enumerate}[leftmargin=1.5em]
    \item 관리자 대시보드에서 \menu{📂 백업 복원} 버튼을 클릭합니다.
    \item \menu{백업 파일 선택}에서 이 시스템에서 내보낸 \texttt{.xlsx} 파일을 선택합니다.
    \item 복원할 항목을 선택합니다.
    \begin{itemize}
        \item \textbf{학생 계정}: 학생명단 시트에서 학생 계정 생성
        \item \textbf{교사 계정}: 교사명단 시트에서 교사 계정·승인 상태·담당 학년 복원 (비밀번호 그대로 복원)
        \item \textbf{공휴일}: 공휴일 시트에서 날짜·이름 복원
        \item \textbf{자습 시간 설정}: 자습시간설정 시트에서 교시별 시간 복원 (이미 존재하면 값 업데이트)
        \item \textbf{출결 기록}: 출결상세 시트에서 원시 출결 데이터 복원
        \item \textbf{자습 신청 기록}: 자습신청 시트에서 복원
        \item \textbf{학습 기록}: 학습기록 시트에서 복원
        \item \textbf{자습실 배정 및 좌석 위치}: 자습실배정 시트에서 배정 공간·좌석번호·배치도 위치 복원
    \end{itemize}
    \item \menu{복원 실행} 버튼을 클릭하고 완료될 때까지 기다립니다.
    \item 완료 후 복원 결과(생성 건수, 중복 건너뜀 수)가 표시됩니다.
\end{enumerate}

\subsubsection{복원 규칙}

\begin{itemize}[leftmargin=1.5em]
    \item \textbf{중복 방지}: 같은 학번의 학생·같은 아이디의 교사가 이미 있으면 건너뜁니다. 같은 날짜·교시의 출결·신청도 중복 삽입하지 않습니다.
    \item \textbf{임시 비밀번호 일괄 발급}: 신규로 생성되는 모든 계정(학생·교사·관리자)에는
    \textbf{영문+숫자 조합 랜덤 임시 비밀번호}가 발급됩니다. 예측 가능한 패턴(\texttt{학번+Study1}
    등)은 사용하지 않습니다. 복원이 끝나면 발급된 명단이 Excel 파일
    \texttt{self\_study\_복원결과\_YYYYMMDD\_HHMMSS.xlsx}로 즉시 다운로드되며,
    \textbf{이 파일은 화면이 아닌 엑셀로만 제공됩니다} (평문 비번이 flash·로그에 남지 않음).
    \item \textbf{자습실 QR 토큰}: 자습실 복원 시 \textbf{QR 토큰이 새로 발급}되므로,
    복원 후에는 교사 페이지에서 QR코드를 다시 인쇄해 부착해야 합니다.
    옛 QR코드는 더 이상 작동하지 않습니다.
    \item \textbf{자습실 복원 순서}: 자습실배정 시트를 복원하려면 자습실 목록이 먼저 복원되어 있어야 합니다. \textbf{자습실 목록}을 함께 선택하면 자동으로 먼저 처리됩니다.
    \item \textbf{순서}: 출결·신청·학습기록·자습실배정을 복원하려면 해당 학생 계정이 먼저 존재해야 합니다. 학생 계정 복원과 함께 선택하거나, 계정 복원 후 다시 실행하십시오.
    \item 같은 파일을 두 번 업로드해도 중복 생성되지 않으므로 안전합니다 (다만 두 번째 업로드에서는 신규 계정이 없어 임시 비번 Excel은 다운로드되지 않습니다).
\end{itemize}

\begin{warnbox}
  다운로드된 \texttt{복원결과} Excel은 모든 신규 계정의 초기 비밀번호를
  \textbf{평문}으로 담고 있습니다. 학생·교사에게 개별 전달한 후
  즉시 삭제하십시오. 이 파일이 외부로 유출되면 해당 계정 전부가 침해됩니다.
\end{warnbox}

\begin{warnbox}
  이 시스템에서 내보낸 백업 파일만 복원을 지원합니다.
  시트 이름이나 열 순서를 임의로 변경한 경우 복원이 실패할 수 있습니다.
\end{warnbox}

\subsection{시스템 설정 (정책값 관리)}

학사 운영 중 자주 조정될 수 있는 정책값을 코드 수정 없이
관리자 화면에서 직접 변경할 수 있는 페이지입니다.
\menu{관리자 대시보드 → 시스템 설정} (\texttt{/admin/system-settings})으로 진입합니다.

\subsubsection*{관리 가능한 항목}

\begin{center}
\begin{tabularx}{\textwidth}{lXcc}
\toprule
\textbf{설정 키} & \textbf{설명} & \textbf{기본값} & \textbf{허용 범위} \\
\midrule
\texttt{early\_checkin\_minutes}      & 사전 입실 허용 시간(분) --- 교시 시작 N분 전부터 입실 QR 허용 & 30   & 0\textendash 120 \\
\texttt{checkout\_grace\_minutes}     & 퇴실 grace 시간(분) --- 교시 종료 후 N분까지 퇴실 QR 허용     & 10   & 0\textendash 60  \\
\texttt{late\_threshold\_minutes}     & 지각 판정 시간(분) --- 교시 시작 후 N분 이내는 출석, 이후는 지각 & 10   & 0\textendash 60  \\
\texttt{apply\_cutoff\_day}           & 월별 자습 신청 마감일 --- 매월 N일 이전이면 다음달 신청 가능     & 20   & 1\textendash 31  \\
\texttt{participation\_rate\_default} & 참여율 통계 기본 기준(\%) --- 통계 화면 첫 진입 시 기본 필터    & 80   & 0\textendash 100 \\
\texttt{password\_min\_length}        & 비밀번호 최소 길이                                            & 8    & 4\textendash 30  \\
\texttt{password\_require\_mixed}     & 비밀번호에 영문+숫자 혼합 강제 여부                            & true & true / false     \\
\texttt{temp\_password\_length}       & 관리자가 비밀번호 초기화 시 자동 생성하는 임시 비번 길이        & 8    & 6\textendash 20  \\
\bottomrule
\end{tabularx}
\end{center}

\subsubsection*{변경 방법}

\begin{enumerate}[leftmargin=1.5em]
    \item \menu{관리자 대시보드 → 시스템 설정}으로 진입합니다.
    \item 변경할 항목의 입력란에 새 값을 입력합니다 (정수형은 허용 범위 내, true/false는 토글).
    \item \menu{변경사항 저장} 버튼을 클릭합니다.
    \item 화면 우측의 \textbf{최근 변경} 컬럼에 변경 시각과 처리한 관리자 이름이 기록됩니다.
\end{enumerate}

\begin{notebox}
  변경한 값은 \textbf{저장 즉시 반영}됩니다 --- 서버 재시작이 필요하지 않습니다.
  예를 들어 \texttt{early\_checkin\_minutes}를 30에서 45로 변경하면,
  바로 다음 학생의 입실 QR 시도부터 45분 전 입실이 허용됩니다.
\end{notebox}

\begin{warnbox}
  \texttt{password\_min\_length}나 \texttt{password\_require\_mixed}를 \textbf{더 엄격하게} 변경하더라도
  기존에 저장된 비밀번호는 영향을 받지 않습니다.
  변경 시점 이후 \textbf{새로 설정·변경되는 비밀번호}부터 새 정책이 적용됩니다.
\end{warnbox}

\subsubsection*{초기 적용}

기존 운영 DB에 \texttt{system\_settings} 테이블이 없는 경우,
\texttt{migrate\_add\_settings.py}를 한 번 실행하여 테이블을 생성하고 기본값을 시드해야 합니다
(자세한 내용은 \texttt{데이터베이스 관리 → DB 마이그레이션} 절 참조).
신규 설치이거나 앱이 한 번 이상 부팅된 적이 있다면 자동으로 시드되므로 별도 작업은 필요하지 않습니다.

\subsection{관리자 마이페이지 (비밀번호 변경)}

\begin{enumerate}[leftmargin=1.5em]
    \item 상단 네비게이션에서 \menu{마이페이지}를 클릭합니다.
    \item \textbf{현재 비밀번호}, \textbf{새 비밀번호}, \textbf{새 비밀번호 확인}을 입력합니다.
    \item 비밀번호 규칙: \texttt{password\_min\_length}자 이상, 영문 + 숫자 포함
          (정확한 길이·혼합 여부는 \menu{시스템 설정}에서 변경 가능).
    \item \menu{변경하기} 버튼을 클릭합니다.
\end{enumerate}

\subsection{관리자 비밀번호 분실 시 재설정}

관리자 비밀번호를 잊어버린 경우, 서버 PC에서 터미널을 열고 다음을 실행합니다.

\begin{lstlisting}[language=bash, caption=비밀번호 직접 재설정]
cd self_study
python -c "
from app import create_app
from models import db, User
app = create_app()
with app.app_context():
    admin = User.query.filter_by(role='admin').first()
    admin.set_password('새비밀번호123')
    db.session.commit()
    print('비밀번호가 변경되었습니다.')
"
\end{lstlisting}

%========================================================
\section{교사 기능}
%========================================================

\subsection{교사 회원가입}

\begin{enumerate}[leftmargin=1.5em]
    \item 로그인 화면에서 \menu{회원가입} 링크를 클릭합니다.
    \item 아이디, 비밀번호, 이름을 입력하고 \textbf{역할}을 \menu{교사}로 선택합니다.
    \item \textbf{담당 학년}을 선택합니다 (선택 사항).
    \item \menu{회원가입} 버튼을 클릭합니다.
    \item ``관리자 승인 후 로그인할 수 있습니다''라는 안내 메시지가 표시됩니다.
    \item 관리자에게 계정 승인을 요청합니다.
    \item 승인 완료 후 로그인이 가능합니다.
\end{enumerate}

\begin{notebox}
  교사 계정은 관리자가 승인하기 전까지 로그인이 불가능합니다.
  미승인 상태에서 로그인을 시도하면 \textbf{승인 대기 안내 화면}이 표시됩니다.
\end{notebox}

\begin{tipbox}
  \textbf{담당 학년} 설정 시, 학생 관리 페이지에 접속하면 해당 학년 학생들만 자동으로 필터링됩니다.
  담당 학년을 선택하지 않으면 전체 학년 학생을 관리할 수 있습니다.
  학년 필터는 페이지에서 직접 변경할 수 있습니다.
\end{tipbox}

\subsection{교사 대시보드}

로그인 후 자동으로 교사 대시보드로 이동합니다. 다음 정보가 표시됩니다.

\begin{itemize}[leftmargin=1.5em]
    \item 전체 학생 수, 오늘 출석 학생 수
    \item 학년별 학생 수
    \item 교시별 오늘 출석 현황
\end{itemize}

\subsection{자습 시간 설정}

자습 교시별 시작·종료 시간을 요일 단위로 자유롭게 설정합니다.

\begin{enumerate}[leftmargin=1.5em]
    \item 상단 네비게이션에서 \menu{설정}을 클릭합니다.
    \item 탭에서 설정할 요일 또는 유형을 선택합니다.\\
          \textbf{월 / 화 / 수 / 목 / 금} — 해당 요일 개별 설정\\
          \textbf{평일(공통)} — 개별 설정이 없는 모든 평일에 적용\\
          \textbf{토요일 / 공휴일} — 각각 별도 설정
    \item 각 교시의 \textbf{활성화} 체크박스, \textbf{시작 시간}, \textbf{종료 시간}을 입력합니다.
    \item (선택) \textbf{동일하게 적용할 요일} 체크박스에서 같은 설정을 쓸 다른 요일을 선택합니다.
    \item 저장 버튼을 클릭합니다.
\end{enumerate}

\subsubsection{요일별 설정 우선순위}

\begin{center}
\begin{tabular}{ll}
\toprule
\textbf{우선순위} & \textbf{적용 조건} \\
\midrule
1순위 & 해당 요일 개별 설정(월\textasciitilde금 탭에서 직접 저장한 경우) \\
2순위 & 평일(공통) 설정 \\
3순위 & 시스템 기본값 \\
\bottomrule
\end{tabular}
\end{center}

예를 들어 수요일에 수업이 일찍 끝나 자습이 15:40에 시작한다면,
\menu{수} 탭에서 교시 시간을 설정하고 저장합니다.
나머지 평일(월·화·목·금)에는 \menu{평일(공통)} 설정이 그대로 적용됩니다.

\subsubsection{여러 요일에 동일 설정 적용}

시간표가 같은 요일이 여러 개일 때 각 탭을 따로 저장하는 번거로움을 줄이기 위해
\textbf{동일하게 적용할 요일} 체크박스를 사용합니다.

\begin{enumerate}[leftmargin=1.5em]
    \item 원하는 요일 탭(예: \menu{월})에서 시간을 설정합니다.
    \item 하단 \textbf{동일하게 적용할 요일} 영역에서 같은 시간을 쓸 요일을 체크합니다\\
          (예: 화, 목, 금, 평일(공통) 체크).
    \item 저장 버튼을 클릭하면 선택한 모든 요일에 동일하게 저장됩니다.
\end{enumerate}

\subsubsection{개별 설정 삭제 (평일 공통으로 되돌리기)}

특정 요일의 개별 설정이 더 이상 필요 없으면 해당 탭 하단의
\menu{개별 설정 삭제 (공통으로 되돌리기)} 버튼을 클릭합니다.
개별 설정이 삭제되면 그 요일은 다시 \textbf{평일(공통)} 설정을 따릅니다.

\subsubsection{요일별 설정 현황 확인}

설정 페이지 하단의 \textbf{요일별 설정 현황} 표에서
각 요일이 개별 설정을 사용하는지, 평일(공통)을 사용하는지 한눈에 확인할 수 있습니다.
\badge{success}{개별 설정} 배지는 별도 설정이 저장된 요일,
\badge{secondary}{평일(공통)} 배지는 공통 설정을 따르는 요일을 나타냅니다.

\subsubsection{기본 자습 시간표}

\begin{center}
\begin{tabular}{clcc}
\toprule
\multicolumn{2}{c}{\textbf{평일(공통) 기본값}} & \multicolumn{2}{c}{\textbf{토요일 기본값}} \\
\cmidrule(lr){1-2}\cmidrule(lr){3-4}
교시 & 시간 & 교시 & 시간 \\
\midrule
아침 자습 & 07:30 -- 08:30 & 1교시 & 09:00 -- 10:00 \\
1교시     & 18:00 -- 19:00 & 2교시 & 10:10 -- 11:10 \\
2교시     & 19:10 -- 20:10 & 3교시 & 11:20 -- 12:20 \\
3교시     & 20:20 -- 21:20 &       &                 \\
4교시     & 21:30 -- 22:30 &       &                 \\
\bottomrule
\end{tabular}
\end{center}

\begin{notebox}
  기본값은 DB에 설정이 전혀 없을 때만 사용됩니다.
  설정 저장 후에는 저장된 값이 항상 우선 적용됩니다.\\
  공휴일의 기본 자습 시간은 1교시(09:00--10:00), 2교시(10:10--11:10)입니다.
  일요일은 자습이 없는 날로 고정됩니다.
\end{notebox}

\subsection{공휴일 등록}

\begin{enumerate}[leftmargin=1.5em]
    \item \menu{설정} 페이지 오른쪽 패널에서 \textbf{공휴일 관리} 섹션을 찾습니다.
    \item \textbf{날짜} 입력란에 공휴일 날짜(예: 2025-03-01)를 입력합니다.
    \item \textbf{공휴일 이름} 입력란에 이름(예: 삼일절)을 입력합니다.
    \item \menu{공휴일 등록} 버튼을 클릭합니다.
    \item 등록된 공휴일은 목록에 표시되며, \menu{삭제} 버튼으로 제거할 수 있습니다.
\end{enumerate}

\begin{tipbox}
  공휴일로 등록된 날은 학생 자습 신청 화면에서 빨간색 배지(\badge{danger}{공휴일 이름})로 표시됩니다.
  공휴일 자습 시간은 \menu{설정 → 공휴일} 탭에서 별도 설정합니다.
\end{tipbox}

\subsection{자습 공간 관리}

자습 공간(독서실, 열람실 등)을 등록하고 학생을 배정합니다.

\subsubsection{자습 공간 추가}

\begin{enumerate}[leftmargin=1.5em]
    \item \menu{설정} 페이지에서 \textbf{자습 공간 관리} 섹션을 찾습니다.
    \item \textbf{공간 이름}(예: 1층 독서실)을 입력합니다.
    \item \textbf{수용 인원}을 입력합니다: 전체 인원, 남학생 수, 여학생 수 순서로 입력합니다.
    \item \menu{공간 추가} 버튼을 클릭합니다.
    \item 공간이 추가되면 자동으로 QR코드용 토큰이 생성됩니다.
\end{enumerate}

\begin{notebox}
  남학생 수와 여학생 수는 자리 랜덤 배치 시 좌석 번호 범위로 사용됩니다.
  예를 들어 남학생 수 20, 여학생 수 15로 설정하면,
  남학생은 1\textasciitilde20 중 랜덤 번호, 여학생은 1\textasciitilde15 중 랜덤 번호를 부여받습니다.
\end{notebox}

\subsubsection{공간 활성화/비활성화 및 삭제}

등록된 공간 목록에서 각 행의 버튼으로 관리합니다.

\begin{center}
\begin{tabular}{ll}
\toprule
\textbf{버튼} & \textbf{동작} \\
\midrule
\badge{primary}{QR코드} & QR코드 표시 페이지로 이동 \\
\badge{secondary}{배치도} & 좌석 배치도 페이지로 이동 \\
임의 배정 & 남녀 구분 랜덤 좌석 번호 자동 배정 \\
수동 배정 & 학생별 좌석 번호 직접 지정 \\
\badge{dark}{수정} & 공간 이름·수용 인원(전체/남/여) 수정 \\
\badge{success}{활성화} / \badgedark{warning}{비활성화} & 해당 공간 사용 여부 전환 \\
\badge{danger}{삭제} & 공간 삭제 (배정 기록 삭제, 출결 기록의 공간 참조는 유지) \\
\bottomrule
\end{tabular}
\end{center}

\subsection{학생 자습 공간 배정}

\begin{enumerate}[leftmargin=1.5em]
    \item 상단 네비게이션에서 \menu{학생 관리}를 클릭합니다.
    \item 학년·반·공간 필터를 사용하여 원하는 학생을 찾습니다.
    \item 해당 학생 행의 \menu{공간 배정} 드롭다운에서 자습 공간을 선택합니다.
    \item \menu{배정} 버튼을 클릭합니다.
\end{enumerate}

\begin{notebox}
  배정된 공간 정보는 QR코드 출석 시 자동으로 기록됩니다.
  공간이 배정되지 않은 학생도 자습 신청과 출석은 가능합니다.
\end{notebox}

\subsubsection{담당 학년 필터}

교사가 담당 학년을 설정한 경우, \menu{학생 관리} 페이지에 접속하면 담당 학년 학생들이
자동으로 필터링되어 표시됩니다. 페이지 상단에 \badge{primary}{N학년 담당} 배지가 표시됩니다.
학년 필터 드롭다운으로 다른 학년을 조회할 수 있습니다.

\subsubsection{배정 해제}

잘못 배정된 학생을 미배정 상태로 되돌릴 수 있습니다.

\begin{enumerate}[leftmargin=1.5em]
    \item 해당 학생 행의 \menu{해제} 버튼을 클릭합니다.
    \item 확인 창에서 \key{확인}을 누르면 즉시 배정이 해제됩니다.
    \item 배정이 해제되면 해당 학생의 좌석 번호도 함께 초기화됩니다.
\end{enumerate}

\begin{notebox}
  \menu{해제} 버튼은 이미 자습 공간이 배정된 학생 행에만 표시됩니다.
  미배정 학생에게는 표시되지 않습니다.
\end{notebox}

\subsection{방과후 수업 등록 및 자동 출석}

방과후 수업에 참여하는 학생은 요일·교시를 미리 등록해 두면, 출석 관리 화면에서
버튼 한 번으로 해당 학생 전원을 출석 처리할 수 있습니다.

\subsubsection{방과후 수업 일정 등록}

\begin{enumerate}[leftmargin=1.5em]
    \item 상단 네비게이션에서 \menu{학생 관리}를 클릭합니다.
    \item 학생 목록 우측 \textbf{방과후 수업} 열의 \menu{등록} (또는 \menu{수정}) 버튼을 클릭합니다.
    \item 팝업 창에 요일(월$\sim$금) × 교시 체크박스 그리드가 표시됩니다.
    \item 해당 학생이 참여하는 요일과 교시를 체크합니다.
    \item \menu{저장} 버튼을 클릭하면 즉시 저장됩니다.
\end{enumerate}

\begin{notebox}
  저장 시 기존 일정은 새 일정으로 \textbf{전체 교체}됩니다.
  등록된 일정은 학생 목록의 방과후 수업 열에 ``월/1교시 화/1교시 \ldots'' 형식으로 요약되어 표시됩니다.
\end{notebox}

\begin{warnbox}
  방과후 수업 일정은 교시 설정(\menu{설정} 페이지)에 등록된 교시 목록을 기준으로 표시됩니다.
  새로운 교시(예: 0교시)를 추가했다면 설정 페이지에서 먼저 해당 교시를 활성화하세요.
\end{warnbox}

\subsubsection{방과후 자동출석 처리}

등록된 일정을 바탕으로 해당 날짜·교시에 출석을 일괄 생성합니다.

\begin{enumerate}[leftmargin=1.5em]
    \item 출석 관리 페이지에서 처리할 날짜를 선택합니다.
    \item 우측 상단의 \menu{📚 방과후 자동출석} 버튼을 클릭합니다.
    \item 확인 창에서 \key{확인}을 누르면 즉시 처리됩니다.
\end{enumerate}

처리 규칙은 다음과 같습니다.

\begin{center}
\begin{tabular}{ll}
\toprule
\textbf{조건} & \textbf{처리} \\
\midrule
해당 날짜 요일에 일정이 등록된 학생 & \badge{success}{방과후출결인정} 처리 \\
자습 신청이 없는 경우 & 신청도 자동 생성 \\
이미 \texttt{after\_school}로 처리된 경우 & 건너뜀 (중복 처리 없음) \\
다른 출결 상태(출석·지각·결석 등)가 있는 경우 & 기존 상태와 관계없이 \textbf{방과후출결인정}으로 업데이트 \\
토요일·일요일 & 처리 불가 (안내 메시지 표시) \\
\bottomrule
\end{tabular}
\end{center}

\begin{notebox}
  자동출석으로 처리된 내역은 출결 수정 이력에 ``방과후자동출석''으로 기록됩니다.
  담당 학년이 설정된 교사는 담당 학년 학생만 처리됩니다.
\end{notebox}

\subsection{자습 공간 자리 배정}

학생들이 자습 공간에 배정된 후, 남녀 구분하여 좌석 번호를 배정합니다.
배정 방식은 \textbf{임의 배정}과 \textbf{수동 배정} 두 가지를 선택할 수 있습니다.

\begin{center}
\begin{tabular}{lll}
\toprule
\textbf{방식} & \textbf{버튼} & \textbf{설명} \\
\midrule
임의 배정 & \menu{임의 배정} & 시스템이 좌석 번호를 무작위로 자동 배정 \\
수동 배정 & \menu{수동 배정} & 교사가 각 번호 슬롯에 학생을 직접 지정 \\
\bottomrule
\end{tabular}
\end{center}

\subsubsection{임의 배정}

\begin{enumerate}[leftmargin=1.5em]
    \item \menu{설정} 페이지의 자습 공간 목록에서 \menu{임의 배정} 버튼을 클릭합니다.
    \item 확인 창에서 \key{확인}을 누르면 즉시 랜덤 배치가 실행됩니다.
    \item 남학생·여학생 각각 정원 $n$ 범위 내에서 중복 없이 번호를 배정합니다.
    \item 배정 완료 후 배치도 페이지로 자동 이동합니다.
\end{enumerate}

\begin{itemize}[leftmargin=1.5em]
    \item 남학생과 여학생의 좌석 번호는 \textbf{독립적}으로 관리됩니다 (각각 1번부터 시작)
    \item 다시 실행하면 기존 좌석 번호가 모두 초기화되고 새로 배정됩니다
\end{itemize}

\begin{warnbox}
  학생 수가 정원을 초과하면 경고 메시지가 표시되며, 초과 학생 수만큼 번호 범위가 자동으로 늘어납니다.
  정원에 맞게 학생을 배정한 후 자리 배치를 실행하는 것을 권장합니다.
\end{warnbox}

\subsubsection{수동 배정}

교사가 각 번호 슬롯에 학생을 직접 지정하는 방식입니다.

\begin{enumerate}[leftmargin=1.5em]
    \item \menu{설정} 페이지의 자습 공간 목록에서 \menu{수동 배정} 버튼을 클릭합니다.
    \item 수동 자리 배정 페이지가 열립니다.
    \item 남학생 영역과 여학생 영역이 각각 좌우로 나뉘어 표시됩니다.
    \item 각 번호 슬롯의 드롭다운에서 배정할 학생을 선택합니다.
    \item 같은 학생이 두 자리에 중복 선택되면 빨간 테두리로 표시됩니다.
    \item 미배정 학생은 페이지 하단에 목록으로 표시됩니다.
    \item \menu{💾 배정 저장 후 배치도로 이동} 버튼을 클릭하면 저장 후 배치도 페이지로 이동합니다.
\end{enumerate}

\begin{notebox}
  중복이 있는 경우 저장 시 경고창이 표시됩니다. 중복을 해소한 후 저장하세요.\\
  일부 학생에게만 번호를 배정하고 나머지를 빈 자리로 두는 것도 허용됩니다.
\end{notebox}

\subsection{좌석 배치도 시각화}

임의 배정 또는 수동 배정 후 실제 자습실 내 학생들의 좌석 위치를 시각적으로 표시하고, 드래그 앤 드롭으로 직접 조정할 수 있습니다.

\subsubsection{배치도 열기}

\begin{enumerate}[leftmargin=1.5em]
    \item \menu{설정} 페이지의 자습 공간 목록에서 해당 공간의 \menu{배치도} 버튼을 클릭합니다.
    \item 자습실을 본뜬 격자(grid) 배경 위에 학생별 카드가 표시됩니다.
    \item 위쪽은 \textbf{앞(칠판)} 방향입니다.
\end{enumerate}

\subsubsection{학생 카드 표시 형식}

각 학생 카드에는 다음 정보가 표시됩니다.

\begin{center}
\begin{tabular}{ll}
\toprule
\textbf{항목} & \textbf{내용} \\
\midrule
좌석 번호 & 랜덤 배치된 번호 \\
학번 & 5자리 학번 \\
이름 & 학생 성명 \\
학년/반 & 예: 2-3 \\
\bottomrule
\end{tabular}
\end{center}

남학생 카드는 \textbf{파란색}, 여학생 카드는 \textbf{분홍색}으로 표시됩니다.

\subsubsection{드래그 앤 드롭으로 위치 조정}

\begin{enumerate}[leftmargin=1.5em]
    \item 이동할 학생 카드를 마우스로 클릭하거나 스마트폰에서 터치합니다.
    \item 카드를 원하는 위치로 드래그합니다.
    \item 마우스를 놓으면 카드가 해당 위치에 고정됩니다.
    \item 카드는 배치도 영역 밖으로 벗어날 수 없습니다.
\end{enumerate}

\begin{notebox}
  PC와 스마트폰 모두에서 드래그가 동작합니다. 스마트폰에서는 손가락으로 터치하여 끌면 됩니다.
\end{notebox}

\subsubsection{배치 저장}

\begin{enumerate}[leftmargin=1.5em]
    \item 위치 조정이 완료되면 우측 하단의 \menu{💾 배치 저장} 버튼을 클릭합니다.
    \item 저장이 완료되면 버튼 옆에 \textbf{✓ 저장됨}이 표시됩니다.
\end{enumerate}

\begin{tipbox}
  저장된 배치 정보는 \textbf{모든 교사가 공유}합니다. 한 교사가 위치를 수정하고 저장하면
  다른 교사가 배치도를 열 때도 동일한 배치가 표시됩니다.
\end{tipbox}

\begin{notebox}
  랜덤 배치(\menu{자리 배치} 버튼)를 다시 실행하면 배치도 위치가 자동으로 초기화됩니다.
  새로 배치된 학생들은 자동으로 균등 배치됩니다.
\end{notebox}

\subsubsection{실시간 출결 현황}

배치도 상단의 컨트롤 바에서 날짜와 교시를 선택하면 각 학생 카드에 출결 상태가 실시간으로 표시됩니다.

\textbf{컨트롤 바 구성}

\begin{center}
\begin{tabular}{ll}
\toprule
\textbf{항목} & \textbf{설명} \\
\midrule
날짜 & 출결을 조회할 날짜 (기본값: 오늘) \\
교시 & 조회할 교시 선택. \menu{전체} 선택 시 당일 최선 상태 표시 \\
⟳ 새로고침 & 즉시 재조회 \\
자동 갱신 & 켜짐 시 30초마다 자동으로 재조회 \\
최종 갱신 & 마지막으로 데이터를 가져온 시각 \\
\bottomrule
\end{tabular}
\end{center}

\textbf{카드 색상 기준}

\begin{center}
\begin{tabular}{lll}
\toprule
\textbf{테두리 색} & \textbf{상태} & \textbf{설명} \\
\midrule
\textcolor{success}{초록} & 출석 & QR 또는 교사가 출석 처리 \\
\textcolor{warning!70!black}{주황} & 지각 & 지각 처리됨 \\
\textcolor{danger}{빨강} & 결석 & 결석 처리됨 \\
\textcolor{primary}{파랑 점선} & 신청(미처리) & 신청했지만 아직 출결 기록 없음 \\
\textcolor{purple}{보라} & 방과후출결인정 & 방과후 수업 참여로 출결 처리됨 \\
회색 & 해당없음 & 해당 날짜·교시에 신청 없음 \\
\bottomrule
\end{tabular}
\end{center}

카드 하단에 상태 색상 점(\textbullet)과 한글 상태명(출석/지각/결석/신청/-)이 표시됩니다.
마우스를 카드 위에 올리면 교시별 상세 상태가 툴팁으로 나타납니다
(예: ``1교시: 출석 / 2교시: 지각 / 3교시: 결석'').

\textbf{전체 교시 모드}

교시를 \menu{전체}로 설정하면 해당 학생의 당일 모든 교시 중
가장 좋은 상태(출석 $>$ 지각 $>$ 신청 $>$ 결석 $>$ 없음)를 대표 색상으로 표시합니다.

\begin{tipbox}
  자동 갱신이 켜진 상태에서 배치도를 열어 두면 QR 출석이 들어올 때마다
  최대 30초 이내에 해당 학생 카드가 자동으로 \textcolor{success}{초록}으로 바뀝니다.
\end{tipbox}

\subsection{자습 신청 현황 조회}

\begin{enumerate}[leftmargin=1.5em]
    \item 상단 네비게이션에서 \menu{신청 현황}을 클릭합니다.
    \item 상단에서 조회할 \textbf{날짜}를 선택합니다.
    \item 학년·공간 필터로 원하는 그룹만 조회합니다.
    \item 각 학생별 교시별 신청 여부를 확인합니다.
    \item 페이지 상단에 교시별 신청 인원 통계가 표시됩니다.
\end{enumerate}

\subsection{출석 관리}

\begin{enumerate}[leftmargin=1.5em]
    \item 상단 네비게이션에서 \menu{출석 관리}를 클릭합니다.
    \item 조회할 날짜를 선택합니다.
    \item 필요 시 \textbf{토요일}/\textbf{공휴일} 체크박스로 일자 유형을 수동 지정합니다.
    \item 출석 상태를 변경하려면 해당 셀의 드롭다운에서 상태를 선택하면 즉시 저장됩니다.
\end{enumerate}

\subsubsection{출석 상태 종류}

\begin{center}
\begin{tabular}{lll}
\toprule
\textbf{상태} & \textbf{코드} & \textbf{설명} \\
\midrule
출석     & \texttt{present}        & 정시에 입실 QR 스캔 또는 교사 직접 기록 \\
지각     & \texttt{late}           & 자습 시간 이후 도착 \\
결석     & \texttt{absent}         & 신청했으나 출석하지 않은 경우 \\
조퇴     & \texttt{early\_leave}   & 자습 종료 전 퇴실 QR 스캔으로 자동 처리 \\
출석인정 & \texttt{approved\_leave} & 교사가 조퇴 사유를 확인하여 출석으로 인정 \\
방과후출결인정 & \texttt{after\_school} & 방과후 수업 참여로 자동 출석 처리 (보라색 표시) \\
\bottomrule
\end{tabular}
\end{center}

\begin{tipbox}
  QR코드로 자동 처리된 출석은 \menu{출석} 상태로 기록됩니다.
  자습 종료 전 퇴실 QR을 스캔하면 \menu{조퇴}로 기록되며,
  교사가 출석 관리 페이지에서 사유를 입력하고 \menu{출석인정}으로 변경할 수 있습니다.
\end{tipbox}

\subsubsection{출석 현황 요약}

출석 관리 페이지 상단에 오늘 날짜 기준 출석 현황을 한눈에 확인할 수 있는 요약 카드가 표시됩니다.

\begin{center}
\begin{tabular}{ll}
\toprule
\textbf{카드} & \textbf{내용} \\
\midrule
전체 요약 & 오늘 출석 인원 / 전체 학생 수 \\
학년별 카드 & 각 학년의 출석 인원 / 전체 인원 \\
자습실별 카드 & 각 자습실의 출석 인원 / 정원 \\
\bottomrule
\end{tabular}
\end{center}

출석 인원이 1명 이상인 카드는 초록(학년별) 또는 파랑(자습실별)으로 강조 표시됩니다.

\subsubsection{출석 관리 내 좌석 배치 현황}

출석 관리 페이지 상단에 자습실별 좌석 배치 현황을 바로 확인할 수 있는 임베드 패널이 제공됩니다.

\begin{enumerate}[leftmargin=1.5em]
    \item 패널 상단의 자습실 탭을 클릭하여 확인할 자습실을 선택합니다.
    \item 교시 선택 드롭다운에서 조회할 교시를 선택합니다.
    \item 학생 카드에 현재 출결 상태가 색상으로 표시됩니다.
    \item \menu{접기/펼치기} 버튼으로 패널을 숨기거나 표시할 수 있습니다.
\end{enumerate}

\begin{tipbox}
  패널은 30초마다 자동으로 갱신됩니다. 출석 관리 페이지를 열어 두면 QR 출석 현황이
  실시간으로 반영됩니다.
\end{tipbox}

\begin{notebox}
  임베드 패널의 카드 색상 기준은 \textbf{좌석 배치도 시각화} 절의 실시간 출결 현황과 동일합니다.
\end{notebox}

\subsubsection{자동 지각/결석 처리}

신청 학생 중 아직 출석 처리가 되지 않은 학생을 시간 기준으로 일괄 처리합니다.

\begin{enumerate}[leftmargin=1.5em]
    \item 출석 관리 페이지 우측 상단의 \menu{⚡ 자동처리} 버튼을 클릭합니다.
    \item 확인 창에서 \key{확인}을 누르면 즉시 처리됩니다.
\end{enumerate}

자동 처리 기준은 다음과 같습니다.

\begin{center}
\begin{tabular}{ll}
\toprule
\textbf{상황} & \textbf{처리 결과} \\
\midrule
오늘 날짜 + 교시 시작 후 10분 이내 미출석 & \textbf{지각} 처리 \\
오늘 날짜 + 교시 종료 후 미출석 & \textbf{결석} 처리 \\
과거 날짜(오늘 이전) & 해당 교시 신청자 모두 \textbf{결석} 처리 \\
\bottomrule
\end{tabular}
\end{center}

\begin{notebox}
  이미 출석·지각·결석 처리된 학생은 자동처리에서 제외됩니다.
  자동처리로 변경된 내역은 모두 \textbf{출결 수정 이력}에 ``자동처리''로 기록됩니다.
\end{notebox}

\subsubsection{방과후 자동출석}

방과후 수업 일정이 등록된 학생을 해당 날짜·교시에 일괄 출석 처리합니다.
자세한 내용은 \textbf{방과후 수업 등록 및 자동 출석} 절을 참고하세요.

\begin{enumerate}[leftmargin=1.5em]
    \item 출석 관리 페이지 우측 상단의 \menu{📚 방과후 자동출석} 버튼을 클릭합니다.
    \item 확인 창에서 \key{확인}을 누르면 즉시 처리됩니다.
\end{enumerate}

\subsubsection{출결 수정 이력}

출결 상태가 변경될 때마다 이력이 자동으로 기록됩니다.

\begin{enumerate}[leftmargin=1.5em]
    \item 출석 관리 페이지 하단의 \textbf{출결 수정 이력} 카드를 찾습니다.
    \item \menu{펼치기} 버튼을 클릭하면 해당 날짜의 이력 목록이 펼쳐집니다.
\end{enumerate}

이력 테이블에 표시되는 항목은 다음과 같습니다.

\begin{center}
\begin{tabularx}{\textwidth}{lX}
\toprule
\textbf{열} & \textbf{내용} \\
\midrule
시각 & 변경이 발생한 시각 (HH:MM:SS) \\
학생 & 출결이 변경된 학생 이름 \\
교시 & 변경된 교시 번호 \\
변경 전 & 변경 이전 상태 (없으면 -- 표시) \\
변경 후 & 변경된 상태 \\
처리자 & 변경한 교사 이름 (자동처리의 경우 시스템) \\
비고 & ``자동처리'' 또는 ``수동처리'' \\
\bottomrule
\end{tabularx}
\end{center}

\subsubsection{출석 현황 Excel 내보내기}

\begin{enumerate}[leftmargin=1.5em]
    \item 출석 관리 페이지 우측 상단의 \menu{📥 Excel} 버튼을 클릭합니다.
    \item \texttt{attendance\_YYYY-MM-DD.xlsx} 파일이 자동으로 다운로드됩니다.
\end{enumerate}

Excel 파일 구성은 다음과 같습니다.

\begin{center}
\begin{tabular}{ll}
\toprule
\textbf{열} & \textbf{내용} \\
\midrule
이름, 학년/반, 자습 공간 & 학생 정보 \\
교시별 출석 상태 & 출석(초록)/지각(노랑)/결석(빨강) 색상 구분 \\
\bottomrule
\end{tabular}
\end{center}

\subsubsection{기간별 전체 출결 Excel 내보내기}

날짜 범위와 교시를 직접 지정하여 전체 학생의 출결 현황을 한 파일로 내보냅니다.
자습실 구분 없이 모든 학생이 포함됩니다.

\begin{enumerate}[leftmargin=1.5em]
    \item 출석 관리 페이지 우측 상단의 \menu{📊 기간별 출결 Excel} 버튼을 클릭합니다.
    \item 팝업 창에서 \textbf{시작일}과 \textbf{종료일}을 입력합니다.
    \item 내보낼 \textbf{교시}를 체크박스로 선택합니다 (기본: 전체 선택).
    \item \menu{📥 Excel 다운로드} 버튼을 클릭하면 파일이 자동으로 다운로드됩니다.
\end{enumerate}

\textbf{Excel 파일 행 구조}

\begin{center}
\begin{tabular}{ll}
\toprule
\textbf{행} & \textbf{내용} \\
\midrule
1행 & 제목 (기간 및 선택 교시 표시) \\
2행 & 날짜 병합 헤더 (예: \texttt{04/15 (화)}) \\
3행 & 교시 헤더 (예: \texttt{1교시}, \texttt{2교시}) \\
4행 & \textbf{감독교사 서명} 칸 --- 고정열 병합 레이블 + 교시별 서명 공간 (연노랑 배경) \\
5행$\sim$ & 학생별 출결 데이터 \\
\bottomrule
\end{tabular}
\end{center}

\textbf{열 구성}

\begin{center}
\begin{tabular}{ll}
\toprule
\textbf{열} & \textbf{내용} \\
\midrule
A (이름) & 학생 성명 \\
B (학번) & 5자리 학번 \\
C (학년) & 학년 \\
D (반) & 반 번호 \\
E (자습공간) & 배정된 자습실 (미배정이면 ``미배정'') \\
F열 이후 & 날짜 × 교시별 출결 상태 (색상 구분) \\
\bottomrule
\end{tabular}
\end{center}

출결 상태별 셀 색상은 다음과 같습니다.

\begin{center}
\begin{tabular}{ll}
\toprule
\textbf{상태} & \textbf{색상} \\
\midrule
출석 & 연초록 \\
지각 & 연노랑 \\
결석 & 연빨강 \\
조퇴 & 연주황 \\
출석인정 & 연파랑 \\
방과후출결인정 & 연보라 \\
기록 없음 & \texttt{-} (색상 없음) \\
\bottomrule
\end{tabular}
\end{center}

날짜 경계마다 굵은 세로선이 표시되어 날짜 구분이 명확합니다.
파일명 형식: \texttt{출결현황\_YYYYMMDD\_YYYYMMDD.xlsx}

\begin{notebox}
  교시를 선택하지 않으면 다운로드가 실행되지 않습니다.
  기간이 길수록 파일 크기가 커지며 생성에 시간이 걸릴 수 있습니다.
\end{notebox}

\subsection{학생 개인 리포트}

학생별 월별 자습 참여 현황을 상세하게 조회합니다.

\begin{enumerate}[leftmargin=1.5em]
    \item \menu{학생 관리} 페이지의 학생 목록에서 해당 학생 행의 \menu{리포트} 버튼을 클릭합니다.
    \item 기본으로 현재 월의 리포트가 표시됩니다.
    \item \menu{◀} / \menu{▶} 버튼으로 이전/다음 달로 이동할 수 있습니다.
\end{enumerate}

\subsubsection{리포트 구성}

\begin{center}
\begin{tabularx}{\textwidth}{lX}
\toprule
\textbf{항목} & \textbf{내용} \\
\midrule
신청 횟수 & 해당 월 총 자습 신청 교시 수 \\
출석 횟수 & 출석 상태(\texttt{present})인 교시 수 \\
지각 횟수 & 지각 상태(\texttt{late})인 교시 수 \\
결석 횟수 & 결석 상태(\texttt{absent})인 교시 수 \\
참여율 & (출석 + 지각) / 신청 × 100\% \\
총 학습 시간 & 해당 월 학습 기록의 합산 시간 \\
\bottomrule
\end{tabularx}
\end{center}

참여율은 색상으로 구분됩니다.

\begin{center}
\begin{tabular}{ll}
\toprule
\textbf{색상} & \textbf{기준} \\
\midrule
초록 & 90\% 이상 \\
파랑 & 80\% 이상 \\
노랑 & 60\% 이상 \\
빨강 & 60\% 미만 \\
\bottomrule
\end{tabular}
\end{center}

리포트 하단에는 날짜별 출결 상세 테이블과 학습 기록 테이블이 표시됩니다.

\begin{notebox}
  학생 개인 리포트는 교사만 열람할 수 있으며, 학생 본인은 대시보드에서 기본 정보를 확인합니다.
\end{notebox}

\subsection{참여 통계}

월별 자습 참여 우수자를 조회합니다.

\begin{enumerate}[leftmargin=1.5em]
    \item 상단 네비게이션에서 \menu{참여 통계}를 클릭합니다.
    \item 조회할 연도와 월을 선택합니다.
    \item \textbf{최소 참여율}(기본 80\%)을 설정합니다.
    \item 학년 필터를 사용하여 특정 학년만 조회합니다.
\end{enumerate}

표에서 학생별로 \textbf{신청 횟수}, \textbf{출석 횟수}, \textbf{참여율},
\textbf{총 자습 시간}을 확인할 수 있으며,
참여율 기준 이상인 학생만 목록에 포함됩니다.

\subsubsection{총 자습 시간}

\textbf{총 자습 시간}은 학생이 입실 QR과 퇴실 QR을 \textbf{모두} 스캔한 교시의
실제 체류 시간을 합산한 값입니다.

\begin{center}
\begin{tabular}{ll}
\toprule
\textbf{조건} & \textbf{처리} \\
\midrule
입실·퇴실 QR 모두 스캔 & 퇴실 시각 $-$ 입실 시각 (분 단위) 누적 \\
퇴실 QR 미스캔 & 해당 교시 제외 (시간 불명) \\
조퇴·출석인정 포함 & 실제 체류 시간이므로 포함 \\
\bottomrule
\end{tabular}
\end{center}

\begin{notebox}
  퇴실 QR을 스캔하지 않은 교시는 총 자습 시간 집계에서 제외됩니다.
  퇴실 QR 스캔을 생활화하도록 학생들에게 안내하십시오.
\end{notebox}

\subsubsection{참여 통계 Excel 내보내기}

\begin{enumerate}[leftmargin=1.5em]
    \item 참여 통계 결과 테이블 아래의 \menu{📥 Excel 다운로드} 버튼을 클릭합니다.
    \item \texttt{참여통계\_YYYY년MM월.xlsx} 파일이 자동으로 다운로드됩니다.
\end{enumerate}

Excel 파일에는 순위, 이름, 학번, 학년, 반, 신청 횟수, 출석 횟수, 참여율,
\textbf{총 자습 시간}이 포함되며,
참여율에 따라 셀 색상이 자동으로 적용됩니다(90\%↑ 초록, 80\%↑ 파랑, 그 외 노랑).

\begin{tipbox}
  \menu{클립보드에 복사} 버튼을 사용하면 Excel에 직접 붙여넣기(Ctrl+V)할 수 있는 형식으로
  데이터가 클립보드에 저장됩니다. 총 자습 시간도 함께 복사됩니다.
\end{tipbox}

\subsection{마이페이지 (교사)}

교사 본인의 계정 정보를 확인하고 비밀번호를 변경합니다.

\begin{enumerate}[leftmargin=1.5em]
    \item 상단 네비게이션에서 \menu{마이페이지}를 클릭합니다.
    \item 상단 카드에서 이름, 아이디, 담당 학년을 확인합니다.
    \item 비밀번호를 변경하려면 \textbf{현재 비밀번호}, \textbf{새 비밀번호},
          \textbf{새 비밀번호 확인}을 입력한 후 \menu{변경하기}를 클릭합니다.
    \item 비밀번호 규칙: 8자 이상, 영문 + 숫자 포함
\end{enumerate}

\subsection{QR코드 생성 및 인쇄}

하나의 자습 공간에 \textbf{입실 QR}과 \textbf{퇴실 QR} 두 종류가 생성됩니다.
같은 토큰을 공유하되 URL 경로가 다릅니다.

\begin{center}
\begin{tabular}{lll}
\toprule
\textbf{종류} & \textbf{색상} & \textbf{URL 경로} \\
\midrule
입실 QR & 파란색 & \texttt{/student/qr-attend/[토큰]} \\
퇴실 QR & 주황색 & \texttt{/student/qr-checkout/[토큰]} \\
\bottomrule
\end{tabular}
\end{center}

\begin{enumerate}[leftmargin=1.5em]
    \item \menu{설정} 페이지의 자습 공간 목록에서 해당 공간의 \menu{QR코드} 버튼을 클릭합니다.
    \item 입실 QR(파란색)과 퇴실 QR(주황색)이 나란히 표시됩니다.
    \item \menu{인쇄하기} 버튼을 클릭하여 인쇄합니다.
    \item 인쇄된 QR코드를 자습실 입구에 게시합니다.
    \begin{itemize}
        \item \textbf{입실 QR}: 자습실 입구 진입 방향 기준 왼쪽(또는 잘 보이는 곳)
        \item \textbf{퇴실 QR}: 입구 출구 방향 기준 또는 오른쪽
    \end{itemize}
\end{enumerate}

\begin{notebox}
  QR코드를 재생성하면 \textbf{입실·퇴실 QR코드 모두 즉시 무효화}됩니다.
  재생성 후에는 두 QR코드를 모두 다시 인쇄하여 교체해야 합니다.
\end{notebox}

%========================================================
\section{학생 기능}
%========================================================

\subsection{학생 회원가입}

\begin{enumerate}[leftmargin=1.5em]
    \item 로그인 화면에서 \menu{회원가입} 링크를 클릭합니다.
    \item 다음 정보를 입력합니다.
    \begin{itemize}
        \item \textbf{아이디}: 영문·숫자 조합 (고유해야 함)
        \item \textbf{비밀번호}: 8자 이상, 영문 + 숫자 포함
        \item \textbf{비밀번호 확인}: 동일하게 재입력
        \item \textbf{이름}: 실명
        \item \textbf{역할}: \menu{학생} 선택
        \item \textbf{학번}: 숫자 5자리 입력
        \item \textbf{성별}: \menu{남} 또는 \menu{여} 선택
        \item \textbf{학년}: 1, 2, 3 중 선택
        \item \textbf{반}: 반 번호 입력
    \end{itemize}
    \item \menu{회원가입} 버튼을 클릭합니다.
    \item 가입 즉시 로그인이 가능합니다.
\end{enumerate}

\begin{notebox}
  \textbf{학번 입력 방법}: 학번은 \textbf{숫자 5자리}만 입력합니다.
  학교에 따라 입학년도가 앞에 붙는 형식(예: \texttt{2024-10101})을 사용하는 경우,
  입학년도를 제외한 뒤 5자리(\texttt{10101})만 입력합니다.
  학번은 동명이인을 구별하는 데 사용되므로 정확하게 입력하십시오.
  이미 등록된 학번으로는 가입할 수 없습니다.
\end{notebox}

\subsection{학생 대시보드}

로그인 후 자동으로 학생 대시보드로 이동합니다. 다음 정보가 표시됩니다.

\begin{center}
\begin{tabular}{ll}
\toprule
\textbf{카드} & \textbf{표시 내용} \\
\midrule
오늘 자습 시간 & 오늘의 교시별 자습 시간과 신청 여부 \\
출석 현황 & 교시별 출석 상태 및 출석 시각 \\
이번 달 신청 현황 & 해당 월의 총 자습 신청 건수 \\
이번 주 학습 현황 & 이번 주 총 학습 기록 시간 \\
\bottomrule
\end{tabular}
\end{center}

\subsection{월별 자습 신청}

\begin{enumerate}[leftmargin=1.5em]
    \item 상단 네비게이션 또는 대시보드에서 \menu{자습 신청}을 클릭합니다.
    \item 기본으로 \textbf{다음 달}이 표시됩니다 (20일 이전인 경우).
    \item \menu{이전}/\menu{다음} 버튼으로 월을 이동합니다.
    \item 달력에서 원하는 날짜와 교시의 체크박스를 선택합니다.
    \begin{itemize}
        \item \textbf{전체} 체크박스: 해당 날짜의 모든 교시 일괄 선택/해제
        \item 각 교시 체크박스: 개별 교시 선택
    \end{itemize}
    \item \menu{신청하기} 버튼을 클릭합니다.
\end{enumerate}

\begin{notebox}
  신청을 다시 제출하면 해당 월의 \textbf{기존 신청 내역이 모두 초기화}되고 새로 저장됩니다.
  수정이 필요하면 변경 사항을 모두 체크한 후 다시 신청하면 됩니다.
\end{notebox}

\begin{warnbox}
  일요일 칸은 선택할 수 없습니다. 공휴일은 빨간색 배지로 표시되며,
  공휴일 교시 설정에 따라 자습 신청이 가능합니다.
\end{warnbox}

\subsection{QR코드로 입실하기 (출석)}

\begin{enumerate}[leftmargin=1.5em]
    \item 자습실에 입장합니다.
    \item 입구에 게시된 \textbf{입실 QR코드(파란색)}를 스마트폰 카메라로 스캔합니다.
    \item 로그인하지 않은 경우 로그인 화면이 나타납니다. 로그인합니다.
    \item 현재 자습 시간이 맞으면 자동으로 출석이 처리되고 대시보드로 이동합니다.
    \item 대시보드에서 입실 시각을 확인할 수 있습니다.
\end{enumerate}

\subsubsection{QR 입실 조건}

다음 조건이 \textbf{모두} 충족되어야 출석이 처리됩니다.

\begin{itemize}[leftmargin=1.5em]
    \item 현재 시각이 해당 교시 시작 시각 \textbf{30분 전} 이후이거나 자습 시간 범위 내에 있을 것
    \item 해당 날짜·교시에 자습 신청이 되어 있을 것
    \item 아직 입실 처리되지 않았을 것
\end{itemize}

\begin{notebox}
  \textbf{사전 입실}: 교시 시작 최대 30분 전부터 QR코드를 스캔할 수 있습니다.
  예를 들어 1교시가 18:00에 시작하면 17:30부터 입실 QR이 동작합니다.
  사전 입실한 경우 ``N교시 사전 입실 완료 (시작: HH:MM)'' 안내가 표시됩니다.
\end{notebox}

\begin{tipbox}
  교시 사이 쉬는 시간에 QR을 스캔하면 자동으로 \textbf{다음 신청 교시}로 사전 입실 처리됩니다.
  이미 출석한 경우에는 ``이미 출석 완료''라고 안내됩니다.
  모든 신청 교시가 종료된 이후 스캔하면 ``출석할 수 있는 교시가 없습니다'' 안내가 표시됩니다.
\end{tipbox}

\subsection{QR코드로 퇴실하기 (조퇴 처리)}

자습 도중 부득이하게 퇴실해야 할 경우 퇴실 QR코드를 스캔합니다.

\begin{enumerate}[leftmargin=1.5em]
    \item 퇴실하기 전 입구에 게시된 \textbf{퇴실 QR코드(주황색)}를 스캔합니다.
    \item 서버가 현재 시각과 교시 종료 시각을 비교합니다.
    \begin{itemize}
        \item \textbf{종료 전 퇴실}: 출결 상태가 \menu{조퇴}로 변경되고 퇴실 시각이 기록됩니다.
        \item \textbf{종료 후 퇴실 (10분 이내)}: 퇴실 시각만 기록되며 출석 상태는 그대로 유지됩니다.
    \end{itemize}
    \item 조퇴 처리된 경우 대시보드에 다음 안내가 표시됩니다.\\
          ``\textit{N교시 조퇴 처리되었습니다. 담당 교사가 사유 확인 후 출석인정 가능합니다.}''
\end{enumerate}

\subsubsection{조퇴 처리 조건}

\begin{itemize}[leftmargin=1.5em]
    \item 해당 교시 입실 기록(출석)이 먼저 있어야 합니다.
    \item 퇴실 QR은 교시 시작 시각부터 \textbf{종료 후 10분}까지 동작합니다.\\
          예: 교시 종료가 20:30이면 20:40까지 퇴실 QR 스캔 가능.
    \item 퇴실 QR을 여러 번 스캔하면 \textbf{마지막 스캔 시각}이 최종 퇴실 시각으로 기록됩니다.
\end{itemize}

\begin{warnbox}
  퇴실 QR은 반드시 \textbf{입실 QR 스캔 후}에만 사용할 수 있습니다.
  입실 기록 없이 퇴실 QR을 스캔하면 ``입실 기록이 없습니다''라는 안내가 표시됩니다.
\end{warnbox}

\subsection{학습 기록}

자율학습 중 공부한 내용을 기록합니다.

\begin{enumerate}[leftmargin=1.5em]
    \item 상단 네비게이션에서 \menu{학습 기록}을 클릭합니다.
    \item 다음 항목을 입력합니다.
    \begin{itemize}
        \item \textbf{날짜}: 기록할 날짜 (기본값: 오늘)
        \item \textbf{과목}: 공부한 과목명
        \item \textbf{학습 시간(분)}: 학습한 시간(분 단위)
        \item \textbf{메모}: 학습 내용 요약 (선택)
    \end{itemize}
    \item \menu{기록 저장} 버튼을 클릭합니다.
    \item 페이지 하단에 최근 30일간의 학습 기록이 표시됩니다.
\end{enumerate}

\subsection{내 출결 현황 (월별 자기확인)}

학생이 본인의 월별 출결 이력을 직접 조회할 수 있습니다.

\begin{enumerate}[leftmargin=1.5em]
    \item 상단 네비게이션에서 \menu{내 출결}을 클릭합니다.
    \item 기본으로 이번 달이 표시됩니다.
    \item \menu{◀} / \menu{▶} 버튼으로 이전·다음 달로 이동합니다.
\end{enumerate}

\subsubsection{월간 요약 카드}

페이지 상단에 다음 6개의 요약 카드가 표시됩니다.

\begin{center}
\begin{tabular}{lll}
\toprule
\textbf{카드} & \textbf{색상} & \textbf{내용} \\
\midrule
신청     & 회색  & 해당 월 총 자습 신청 교시 수 \\
출석     & 초록  & 출석(\texttt{present}) 처리된 교시 수 \\
지각     & 노랑  & 지각(\texttt{late}) 처리된 교시 수 \\
결석     & 빨강  & 결석(\texttt{absent}) 처리된 교시 수 \\
조퇴     & 주황  & 조퇴(\texttt{early\_leave}) 처리된 교시 수 \\
출석인정 & 하늘  & 출석인정(\texttt{approved\_leave}) 처리된 교시 수 \\
\bottomrule
\end{tabular}
\end{center}

참여율 = (출석 + 지각 + \textbf{출석인정}) ÷ 신청 × 100\%이며,
비율에 따라 진행 막대 색상이 달라집니다.

\begin{center}
\begin{tabular}{ll}
\toprule
\textbf{색상} & \textbf{기준} \\
\midrule
초록 & 90\% 이상 \\
파랑 & 80\% 이상 \\
노랑 & 60\% 이상 \\
빨강 & 60\% 미만 \\
\bottomrule
\end{tabular}
\end{center}

\subsubsection{일별 출결 상세 테이블}

신청 또는 출결 기록이 있는 날짜가 행으로 나열되고,
교시별 열에 출결 상태 배지가 표시됩니다.

\begin{center}
\begin{tabular}{ll}
\toprule
\textbf{표시} & \textbf{의미} \\
\midrule
\badge{success}{출석}   & 출석 처리됨 (입실 QR 또는 교사 기록) \\
\badgedark{warning}{지각}   & 지각 처리됨 \\
\badge{danger}{결석}    & 결석 처리됨 \\
\badge{orange}{조퇴}    & 자습 종료 전 퇴실 QR 스캔으로 자동 처리 \\
\badgedark{info}{출석인정}  & 교사가 조퇴를 출석으로 인정 \\
\texttt{신청}           & 신청했으나 아직 출결 미처리 \\
\texttt{-}              & 해당 교시 신청 없음 \\
\bottomrule
\end{tabular}
\end{center}

배지 아래에 다음 시각 정보가 표시됩니다.

\begin{itemize}[leftmargin=1.5em]
    \item \textbf{입실 HH:MM}: 입실 QR 스캔 시각 (출석·지각·조퇴·출석인정 공통)
    \item \textbf{퇴실 HH:MM}: 퇴실 QR 스캔 시각 (조퇴 시 주황색으로 표시)
    \item \textbf{📝 사유}: 교사가 출석인정 처리 시 입력한 메모 (출석인정 상태에만 표시)
\end{itemize}

\begin{notebox}
  미래 날짜는 조회 결과에 포함되지 않습니다. 당일까지의 데이터만 표시됩니다.
\end{notebox}

\subsection{마이페이지 (학생)}

학생 본인의 계정 정보를 확인하고 비밀번호를 변경합니다.

\begin{enumerate}[leftmargin=1.5em]
    \item 상단 네비게이션에서 \menu{마이페이지}를 클릭합니다.
    \item 상단 카드에서 이름, 아이디, 학번, 학년/반, 성별을 확인합니다.
    \item 비밀번호를 변경하려면 \textbf{현재 비밀번호}, \textbf{새 비밀번호},
          \textbf{새 비밀번호 확인}을 입력한 후 \menu{변경하기}를 클릭합니다.
    \item 비밀번호 규칙: 8자 이상, 영문 + 숫자 포함
\end{enumerate}

%========================================================
\section{QR코드 입실·퇴실 시스템 상세}
%========================================================

\subsection{시스템 구성}

QR코드 입·퇴실 시스템은 다음과 같이 동작합니다.

\begin{enumerate}[leftmargin=1.5em, label=\arabic*단계:]
    \item 교사가 \menu{설정}에서 자습 공간을 추가하면 고유 토큰이 자동 생성됩니다.
    \item 토큰을 기반으로 두 가지 URL이 생성됩니다.
          \begin{itemize}
              \item 입실: \texttt{/student/qr-attend/[토큰]}
              \item 퇴실: \texttt{/student/qr-checkout/[토큰]}
          \end{itemize}
    \item QR코드 표시 페이지에서 입실·퇴실 QR코드를 함께 인쇄하여 자습실에 게시합니다.
    \item 학생이 입실 QR을 스캔하면 서버가 \textbf{해당 요일의 자습 시간 설정}을 기준으로
          현재 교시를 판별하고 출석을 기록합니다.
          교시 시작 \textbf{30분 전}부터 사전 입실이 허용됩니다.
    \item 학생이 자습 종료 전에 퇴실 QR을 스캔하면 조퇴로 기록되고 퇴실 시각이 저장됩니다.
    \item 교사가 출석 관리 페이지에서 조퇴 학생의 사유를 확인하고 출석인정으로 변경할 수 있습니다.
\end{enumerate}

\subsection{조퇴 및 출석인정 처리 흐름}

\begin{center}
\begin{tabular}{lll}
\toprule
\textbf{단계} & \textbf{처리 주체} & \textbf{결과} \\
\midrule
입실 QR 스캔                  & 학생(자동)    & 출석(\texttt{present}) 기록, 입실 시각 저장 \\
자습 중 퇴실 QR 스캔          & 학생(자동)    & 조퇴(\texttt{early\_leave}) 변경, 퇴실 시각 저장 \\
출결 관리 페이지에서 사유 입력 & 교사          & 사유 텍스트 저장 \\
\menu{출석인정} 버튼 클릭      & 교사          & 출석인정(\texttt{approved\_leave}) 변경, 이력 기록 \\
\bottomrule
\end{tabular}
\end{center}

\begin{notebox}
  조퇴 처리 후 교사가 \menu{출석 관리} 페이지를 열면 해당 학생의 셀에
  ``사유 입력'' 텍스트 필드와 \menu{출석인정} 버튼이 자동으로 표시됩니다.
  사유를 입력하지 않아도 출석인정 처리는 가능합니다.
\end{notebox}

\subsection{QR코드 인쇄 방법}

\begin{enumerate}[leftmargin=1.5em]
    \item \menu{설정} → 자습 공간 목록에서 \menu{QR코드} 버튼 클릭
    \item 좌측 파란색(입실)·우측 주황색(퇴실) QR코드가 나란히 표시됩니다.
    \item \menu{인쇄하기} 버튼 클릭 — 네비게이션 바와 버튼은 자동으로 숨겨집니다.
    \item 브라우저 인쇄 다이얼로그에서 용지 크기를 A4로 설정합니다.
\end{enumerate}

\begin{tipbox}
  QR코드가 작은 경우 스캔이 잘 안 될 수 있습니다.
  A4 용지 절반 크기 이상으로 인쇄하는 것을 권장합니다.
  입실·퇴실 QR코드를 구분하기 위해 색상이 다르게 표시됩니다.
\end{tipbox}

\subsection{네트워크 환경}

\subsubsection{현재 네트워크 구조와 문제}

학교 네트워크는 물리적으로 하나이지만, \textbf{유선 LAN}과 \textbf{학교 관리 Wi-Fi}가
서로 다른 IP 대역(서브넷)으로 분리되어 있습니다.

\begin{center}
\begin{tabular}{lll}
\toprule
\textbf{네트워크} & \textbf{IP 대역 예시} & \textbf{해당 기기} \\
\midrule
서버 유선 LAN & 192.168.0.x & 서버 PC, 서버실 로컬 기기 \\
학교 관리 Wi-Fi & 10.x.x.x & 교사용·학생용 Wi-Fi 접속 기기 \\
\bottomrule
\end{tabular}
\end{center}

두 대역 사이에 라우팅이 없기 때문에, 학교 Wi-Fi에 접속한 기기는
유선 LAN에 있는 서버 PC에 기본적으로 접근할 수 없습니다.

\begin{center}
\ttfamily\small
\begin{tabular}{ccc}
[서버 유선 LAN 192.168.0.x] & $\longleftrightarrow$ & [학교 Wi-Fi 10.x.x.x] \\
서버 PC, 서버실 기기 & \textcolor{danger}{라우팅 없음} & 교사/학생 스마트폰·노트북 \\
\end{tabular}
\end{center}

\subsubsection{서버 관리자가 할 일}

서버 PC를 학교 Wi-Fi에 \textbf{직접 접속}시키면 네트워크 관리자 도움 없이 해결됩니다.
서버 PC가 학교 Wi-Fi에서 IP를 받으면, 학교 Wi-Fi 사용자와 같은 대역이 되어 접속이 가능합니다.
서버는 \texttt{0.0.0.0}으로 바인딩되어 있으므로 유선과 무선 두 인터페이스 모두에서 자동으로 수신합니다.

\begin{enumerate}[leftmargin=1.5em]
  \item \textbf{USB Wi-Fi 동글 준비}\\
    데스크탑 서버 PC는 일반적으로 내장 Wi-Fi 카드가 없으므로
    \textbf{USB Wi-Fi 동글}을 별도 구매하여 서버 PC의 USB 포트에 연결합니다.

    \begin{center}
    \begin{tabular}{ll}
    \toprule
    \textbf{구매 시 확인 사항} & \textbf{내용} \\
    \midrule
    운영체제 지원 & Windows 10/11 드라이버 포함 제품 \\
    규격 & USB 2.0 이상 (USB 3.0 권장) \\
    가격대 & 1만 원 내외 제품으로 충분 \\
    \bottomrule
    \end{tabular}
    \end{center}

    동글 연결 후 드라이버를 설치하면 Windows 작업표시줄에 Wi-Fi 아이콘이 나타납니다.\\
    정상 인식 여부는 관리자 권한 PowerShell에서 확인합니다.
\begin{lstlisting}[language=bash]
Get-NetAdapter | Select-Object Name, InterfaceDescription, Status
\end{lstlisting}
    목록에 \texttt{Wi-Fi} 항목이 보이고 Status가 \texttt{Disconnected}이면 정상 인식된 것입니다.

  \item \textbf{서버 PC를 교사용 Wi-Fi에 접속}\\
    Windows 작업표시줄 우측 하단 Wi-Fi 아이콘 클릭 → 교사용 Wi-Fi 선택 → 비밀번호 입력.\\
    교사용·학생용 Wi-Fi는 같은 서브넷을 사용하므로 교사용에 접속하면 학생도 동일 IP로 접속 가능합니다.

  \item \textbf{방화벽 규칙 추가}\\
    관리자 권한 PowerShell에서 아래 명령을 실행합니다 (최초 1회):
\begin{lstlisting}[language=bash]
New-NetFirewallRule -DisplayName "자율학습 시스템 Flask 5000" `
  -Direction Inbound -Protocol TCP `
  -LocalPort 5000 -Action Allow -Profile Any
\end{lstlisting}

  \item \textbf{서버 실행 및 Wi-Fi IP 확인}\\
    \texttt{python app.py} 실행 후 터미널 출력에서 IP 주소를 확인합니다.\\
    또는 PowerShell에서:
\begin{lstlisting}[language=bash]
ipconfig
# "무선 LAN 어댑터 Wi-Fi"의 IPv4 주소가 학교 Wi-Fi IP
\end{lstlisting}

  \item \textbf{접속 주소 안내}\\
    Wi-Fi IP(예: \texttt{10.x.x.x:5000})를 교사·학생에게 공지합니다.
\end{enumerate}

\begin{tipbox}
  서버 PC에 두 개의 IP가 생깁니다. \\
  \textbf{유선 LAN IP} (192.168.0.x): 서버실 내 유선 연결 기기용 \\
  \textbf{Wi-Fi IP} (10.x.x.x): 학교 교사·학생 Wi-Fi 기기용 \\
  Flask가 두 IP 모두에서 동시에 응답합니다.
\end{tipbox}

\subsubsection{무선 어댑터 자동 끊김 방지}

서버 PC에 유선(인터넷용)과 무선(학교망용) 두 인터페이스가 동시에 연결된 환경에서,
학교 무선 신호가 약해지면 Windows가 무선 어댑터를 자동으로 끊어버리는 현상이 발생할 수 있습니다.

\begin{center}
\begin{tabular}{ll}
\toprule
\textbf{인터페이스} & \textbf{역할} \\
\midrule
유선 LAN (이더넷) & 인터넷 연결 \\
무선 LAN (Wi-Fi) & 학교망 연결 (self\_study 시스템 운영) \\
\bottomrule
\end{tabular}
\end{center}

두 인터페이스가 반드시 \textbf{동시에 연결}되어 있어야 합니다.
Windows가 무선 신호 약화를 감지하면 유선이 있으므로 무선을 절전 모드로 전환하거나
끊어버리는 것이 원인입니다.

\textbf{1단계: 무선 어댑터 전원 관리 비활성화 (가장 중요)}

\begin{enumerate}[leftmargin=1.5em]
  \item \textbf{장치 관리자} → \textbf{네트워크 어댑터} → Wi-Fi 어댑터 우클릭 → \menu{속성}
  \item \textbf{전원 관리} 탭 → ``전원을 절약하기 위해 컴퓨터가 이 장치를 끌 수 있음'' \textbf{체크 해제}
  \item \menu{확인} 클릭
\end{enumerate}

또는 관리자 권한 PowerShell에서 한 번에 적용할 수 있습니다.

\begin{lstlisting}[language=bash]
# Wi-Fi 어댑터 이름 확인
Get-NetAdapter | Select-Object Name, InterfaceDescription, Status

# 전원 관리 비활성화 (어댑터 이름이 "Wi-Fi"인 경우)
Get-NetAdapter -Name "Wi-Fi" |
  Set-NetAdapterPowerManagement -AllowComputerToTurnOffDevice Disabled
\end{lstlisting}

\textbf{2단계: 무선 어댑터 고성능 모드 설정}

장치 관리자 → Wi-Fi 어댑터 → \menu{속성} → \textbf{고급} 탭에서 아래 항목을 설정합니다.

\begin{center}
\begin{tabular}{ll}
\toprule
\textbf{항목 이름} & \textbf{설정값} \\
\midrule
Power Management / Power Saving Mode & Maximum Performance 또는 Disabled \\
Roaming Aggressiveness & Lowest (신호 약해도 현재 AP 유지) \\
Minimum Power Consumption & Disabled \\
\bottomrule
\end{tabular}
\end{center}

\begin{notebox}
  어댑터 제조사마다 항목 이름이 다를 수 있습니다.
  ``Power''가 포함된 항목을 찾아 ``최대 성능'' 또는 ``비활성화''로 설정하십시오.
\end{notebox}

\textbf{3단계: 인터페이스 메트릭 고정 (라우팅 충돌 방지)}

두 인터페이스가 공존하도록 메트릭을 명시적으로 설정합니다.
관리자 권한 명령 프롬프트에서 실행합니다.

\begin{lstlisting}[language=bash]
REM 유선: 인터넷용 (낮은 메트릭 = 기본 경로)
netsh interface ip set interface "이더넷" metric=10

REM 무선: 학교망용 (높은 메트릭 = 인터넷 기본 경로에서 제외)
netsh interface ip set interface "Wi-Fi" metric=20
\end{lstlisting}

\textbf{4단계: Windows 전원 계획 고성능으로 변경}

\begin{lstlisting}[language=bash]
REM 고성능 전원 계획 활성화
powercfg /setactive 8c5e7fda-e8bf-4a96-9a85-a6e23a8c635c
\end{lstlisting}

또는 \texttt{제어판} → \texttt{전원 옵션} → \menu{고성능} 선택.

\textbf{설정 확인}

관리자 권한 PowerShell에서 두 어댑터가 동시에 \texttt{Up} 상태인지 확인합니다.

\begin{lstlisting}[language=bash]
Get-NetAdapter | Where-Object {$_.Status -eq "Up"} |
  Select-Object Name, Status, LinkSpeed
\end{lstlisting}

유선(이더넷)과 무선(Wi-Fi) 모두 \texttt{Up}이면 정상입니다.

\begin{tipbox}
  1단계(전원 관리 비활성화)와 2단계(고성능 모드)가 핵심입니다.
  이 두 단계만 적용해도 대부분의 무선 끊김 현상이 해결됩니다.
\end{tipbox}

\subsubsection{학교 네트워크 관리자가 할 일 (대안)}

서버 PC에 Wi-Fi 어댑터가 없거나, 서버 PC를 Wi-Fi에 접속시키기 어려운 환경이라면
학교 네트워크 담당자에게 다음을 요청합니다.

\begin{notebox}
  \textbf{네트워크 담당자 요청 내용:}\\
  ``학교 Wi-Fi 대역(교사용·학생용)에서 서버 PC의 유선 IP 주소 \texttt{x.x.x.x}의
  \textbf{5000번 TCP 포트}로 접근할 수 있도록 라우팅 또는 방화벽 허용 규칙을 추가해 주세요.''
\end{notebox}

\begin{center}
\begin{tabular}{lp{9cm}}
\toprule
\textbf{방법} & \textbf{설명} \\
\midrule
서버 관리자 자체 해결 & 서버 PC를 학교 Wi-Fi에 직접 접속 (Wi-Fi 어댑터 필요) \\
네트워크 관리자 요청 & 유선 LAN ↔ Wi-Fi 대역 간 5000번 포트 라우팅 허용 요청 \\
\bottomrule
\end{tabular}
\end{center}

%========================================================
\section{보안}
%========================================================

\subsection{CSRF(Cross-Site Request Forgery) 보호}

본 시스템은 \textbf{Flask-WTF} 라이브러리의 \texttt{CSRFProtect}를 사용하여 모든 POST 요청에
위변조 방지 토큰을 적용합니다.

\subsubsection{동작 방식}

\begin{enumerate}[leftmargin=1.5em]
    \item 서버는 각 세션마다 고유한 CSRF 토큰을 생성합니다.
    \item 모든 페이지의 \texttt{<head>}에 다음과 같이 토큰이 삽입됩니다.
\begin{lstlisting}[language=html]
<meta name="csrf-token" content="[토큰값]">
\end{lstlisting}
    \item 페이지에 포함된 JavaScript가 모든 POST 폼 제출 시 자동으로 토큰을 추가합니다.
    \item 좌석 배치도 저장 등 \texttt{fetch()} API 호출 시에도 \texttt{X-CSRFToken} 헤더로 토큰을 전송합니다.
    \item 서버는 토큰이 없거나 불일치하면 \textbf{400 Bad Request}로 거부합니다.
\end{enumerate}

\subsubsection{보호 범위}

CSRF 토큰은 다음 모든 POST 요청에 자동 적용됩니다.

\begin{itemize}[leftmargin=1.5em]
    \item 로그인, 회원가입
    \item 자습 신청, 학습 기록 저장
    \item 출석 수정, 자동처리, QR코드 재생성
    \item 자습 공간 추가·삭제·배정·자리 배치
    \item 공휴일 등록·삭제, 자습 시간 설정
    \item 교사 승인·거부·권한 취소
    \item 사용자 삭제, 비밀번호 초기화·변경
    \item 좌석 배치도 저장 (JSON API)
\end{itemize}

\begin{notebox}
  GET 요청(조회 전용)과 출결 현황 조회 JSON API(\texttt{/teacher/api/\ldots})는
  데이터를 변경하지 않으므로 CSRF 토큰이 필요하지 않습니다.
\end{notebox}

\begin{warnbox}
  브라우저에서 JavaScript가 비활성화된 경우 CSRF 토큰 자동 주입이 동작하지 않아
  폼 제출이 실패합니다. 본 시스템은 JavaScript가 활성화된 환경에서 사용하십시오.
\end{warnbox}

%========================================================
\section{데이터베이스 관리}
%========================================================

\subsection{데이터베이스 파일 위치}

데이터베이스는 SQLite를 사용하며, 파일 위치는 다음과 같습니다.

\begin{lstlisting}
self_study/
└── instance/
    ├── self_study.db      ← 모든 데이터가 저장되는 파일
    └── secret_key.txt     ← 세션 암호화 키 (최초 실행 시 자동 생성)
\end{lstlisting}

\begin{warnbox}
  \texttt{secret\_key.txt}는 로그인 세션을 보호하는 암호화 키입니다.
  이 파일이 삭제되면 서버 재시작 시 새 키가 생성되고 \textbf{모든 사용자가 강제 로그아웃}됩니다.
  데이터베이스 백업 시 함께 보관하십시오.
\end{warnbox}

\subsection{데이터베이스 구조 (주요 테이블)}

\begin{center}
\begin{tabularx}{\textwidth}{lX}
\toprule
\textbf{테이블} & \textbf{저장 내용} \\
\midrule
\texttt{users} & 사용자 정보 (아이디, 비밀번호, 역할, 성별, 승인 여부) \\
\texttt{study\_applications} & 학생별 월별 자습 신청 내역 \\
\texttt{attendance} & 출석 기록 (날짜, 교시, 상태, 공간) \\
\texttt{study\_logs} & 학습 기록 (날짜, 과목, 시간, 메모) \\
\texttt{study\_rooms} & 자습 공간 정보 및 QR 토큰 \\
\texttt{student\_rooms} & 학생별 자습 공간 배정 \\
\texttt{study\_period\_settings} & 평일·토요일·공휴일 자습 시간 설정 \\
\texttt{system\_settings} & 운영 정책값 (사전 입실 시간·지각 임계·신청 마감일·비밀번호 정책 등) \\
\texttt{holidays} & 공휴일 목록 \\
\texttt{attendance\_logs} & 출결 수정 이력 (변경 전후 상태, 처리자, 처리 구분) \\
\bottomrule
\end{tabularx}
\end{center}

\paragraph{주요 데이터 무결성 제약}

v1.9 이상에서 핵심 테이블에 CHECK · UNIQUE 제약이 적용됩니다
(기존 DB는 \texttt{migrate\_add\_constraints\_v2.py}로 소급 적용).

\begin{center}
\begin{tabularx}{\textwidth}{llX}
\toprule
\textbf{테이블} & \textbf{제약} & \textbf{내용} \\
\midrule
\texttt{users}         & CHECK & \texttt{role} $\in$ \{student, teacher, admin\} \\
\texttt{users}         & CHECK & \texttt{gender} $\in$ \{NULL, M, F\} \\
\texttt{attendance}    & CHECK & \texttt{status} $\in$ \{present, late, absent, early\_leave, approved\_leave, after\_school\} \\
\texttt{attendance}    & UNIQUE & \texttt{(user\_id, date, period)} --- 같은 날·교시 중복 방지 \\
\texttt{student\_rooms} & UNIQUE & \texttt{user\_id} --- 한 학생당 한 자습실 \\
\texttt{study\_rooms}  & UNIQUE & \texttt{name} --- 자습실 이름 중복 방지 \\
\texttt{users}         & UNIQUE & \texttt{username}, \texttt{student\_id} \\
\texttt{study\_period\_settings} & UNIQUE & \texttt{(day\_type, period)} \\
\bottomrule
\end{tabularx}
\end{center}

\subsection{기존 데이터베이스 스키마 업그레이드}

v1.0 데이터베이스를 v1.1로 업그레이드하려면 다음 Python 스크립트를 서버 PC 터미널에서 실행합니다.

\begin{lstlisting}[language=bash, caption=DB 스키마 업그레이드]
cd self_study
python -c "
import sqlite3, os
db = os.path.join('instance', 'self_study.db')
con = sqlite3.connect(db)
cur = con.cursor()
# 새 열 추가 (이미 존재하면 무시)
migrations = [
    'ALTER TABLE users ADD COLUMN gender TEXT',
    'ALTER TABLE users ADD COLUMN student_id TEXT',
    'ALTER TABLE users ADD COLUMN assigned_grade INTEGER',
    'ALTER TABLE study_rooms ADD COLUMN male_capacity INTEGER DEFAULT 0',
    'ALTER TABLE study_rooms ADD COLUMN female_capacity INTEGER DEFAULT 0',
    'ALTER TABLE student_rooms ADD COLUMN seat_number INTEGER',
    'ALTER TABLE student_rooms ADD COLUMN pos_x REAL',
    'ALTER TABLE student_rooms ADD COLUMN pos_y REAL',
]
for sql in migrations:
    try: cur.execute(sql)
    except: pass
# attendance_logs 테이블 생성
cur.execute('''CREATE TABLE IF NOT EXISTS attendance_logs (
    id INTEGER PRIMARY KEY,
    attendance_id INTEGER NOT NULL REFERENCES attendance(id) ON DELETE CASCADE,
    changed_by INTEGER REFERENCES users(id),
    old_status TEXT, new_status TEXT NOT NULL,
    changed_at DATETIME, note TEXT)''')
con.commit(); con.close()
print('업그레이드 완료')
"
\end{lstlisting}

\begin{notebox}
  신규 설치 시에는 이 스크립트를 실행할 필요가 없습니다.
  \texttt{reset\_db.py}를 실행하면 최신 스키마로 자동 생성됩니다.
\end{notebox}

\subsection{데이터 초기화}

\begin{warnbox}
  아래 명령은 \textbf{모든 데이터를 삭제}합니다. 반드시 운영 전 또는 테스트 목적으로만 사용하십시오.
\end{warnbox}

\begin{lstlisting}[language=bash]
python reset_db.py
\end{lstlisting}

\subsection{데이터 백업}

정기적으로 데이터베이스 파일을 백업하는 것을 권장합니다.

\begin{lstlisting}[language=bash, caption=백업 예시 (Windows)]
copy self_study\instance\self_study.db self_study.db.bak
\end{lstlisting}

\begin{lstlisting}[language=bash, caption=백업 예시 (macOS/Linux)]
cp self_study/instance/self_study.db self_study.db.bak
\end{lstlisting}

\subsection{배포 패키지 작성 시 제외 파일}

개발 PC에서 운영 PC로 코드를 옮기거나 백업 압축 파일을 만들 때,
다음 파일·폴더는 \textbf{반드시 제외}해야 합니다.
이들이 압축 파일에 들어가면 운영 데이터·인증 정보가 그대로 노출됩니다.

\begin{center}
\begin{tabularx}{\textwidth}{lX}
\toprule
\textbf{경로} & \textbf{제외 이유} \\
\midrule
\texttt{instance/self\_study.db}        & 모든 사용자 비밀번호 해시·QR 토큰·출결·학습 기록 포함 \\
\texttt{instance/self\_study.db-shm}    & SQLite 트랜잭션 메타 (DB 파일과 짝) \\
\texttt{instance/self\_study.db-wal}    & SQLite WAL 로그 (DB 파일과 짝) \\
\texttt{instance/self\_study\_pre\_reset\_*.db} & 학년도 초기화 전 자동 생성된 백업 DB --- 옛 학생 정보 포함 \\
\texttt{instance/secret\_key.txt}        & 세션 암호화 키. 유출 시 세션 위조 가능 \\
\texttt{logs/app.log}                    & 표준 출력 로그 (Waitress) — IP·요청 흐름·예외 메시지 포함 \\
\texttt{logs/audit.log}                  & 구조화 감사 로그 (5MB $\times$ 5개 회전) — 인증·관리자 작업 추적용 \\
\texttt{\_\_pycache\_\_/}                & 파이썬 바이트코드 캐시. 옮겨도 자동 재생성됨 \\
\texttt{*.aux, *.log, *.out, *.toc, *.synctex.gz} & LaTeX 컴파일 부산물 \\
\bottomrule
\end{tabularx}
\end{center}

\subsubsection*{포함해야 할 것}

\texttt{*.py}, \texttt{requirements.txt}, \texttt{templates/}, \texttt{start\_self\_study.bat},
\texttt{manual.tex}, \texttt{manual.pdf}만 압축에 포함합니다.
운영 PC에서는 \texttt{python app.py}를 처음 실행하면
\texttt{instance/self\_study.db}와 \texttt{instance/secret\_key.txt}가 자동 생성됩니다.

\subsubsection*{개발 PC 정기 정리 권장}

개발 PC에서도 다음 산출물은 주기적으로 삭제하는 것이 좋습니다
(재실행 시 자동 재생성되므로 삭제해도 기능에 영향 없음).

\begin{lstlisting}[language=bash, caption=개발 디렉터리 정리 (PowerShell)]
cd C:\python_projects\self_study
Remove-Item -Recurse -Force __pycache__ -ErrorAction SilentlyContinue
Remove-Item manual.aux, manual.log, manual.out, manual.toc, manual.synctex.gz `
            -ErrorAction SilentlyContinue
\end{lstlisting}

\subsubsection*{PowerShell 압축 예시 (제외 파일 처리)}

\begin{lstlisting}[language=bash, caption=Windows PowerShell 안전 압축]
# 임시 폴더에 필요한 파일만 복사 후 압축
$src = "C:\python_projects\self_study"
$tmp = "$env:TEMP\self_study_pkg"
Remove-Item -Recurse -Force $tmp -ErrorAction SilentlyContinue
robocopy $src $tmp /E `
  /XD instance logs __pycache__ `
  /XF *.aux *.log *.out *.toc *.synctex.gz *.zip
Compress-Archive -Path "$tmp\*" -DestinationPath "$src\self_study.zip" -Force
Remove-Item -Recurse -Force $tmp
\end{lstlisting}

\begin{warnbox}
  과거에 위 제외 파일들이 포함된 압축본이 외부(이메일·USB·클라우드)로 전달된 적이 있다면,
  유출을 전제로 다음 작업을 수행해야 합니다.
  \begin{itemize}
    \item \texttt{instance/secret\_key.txt} 삭제 → 다음 부팅 시 새 키 자동 생성
          (모든 사용자가 강제 로그아웃됨)
    \item 관리자 비밀번호 변경
    \item 모든 자습실 QR 토큰 재발급 후 QR 재인쇄
  \end{itemize}
\end{warnbox}

\subsection{관리자 화면에서 백업 및 복원}

시스템 코드 업데이트나 서버 이전 전에는 관리자 대시보드에서 백업을 수행할 수 있습니다.

\subsubsection*{백업 방법}

관리자 대시보드(\texttt{/admin/}) 하단의 \textbf{시스템 업데이트 전 DB 백업} 카드에서
두 가지 방식의 백업을 제공합니다.

\begin{center}
\begin{tabularx}{\textwidth}{lX}
\toprule
\textbf{백업 방식} & \textbf{설명} \\
\midrule
DB 파일 백업 (.db) &
파일명: \texttt{self\_study\_DB\_YYYYMMDD.db} \\
& SQLite DB 전체를 그대로 저장. 모든 테이블·기록 완전 보존. 완전 복원 가능. \\[0.3em]
Excel 백업 (.xlsx) &
파일명: \texttt{self\_study\_backup\_YYYYMMDD.xlsx} \\
& 11개 시트로 구성. 내용 확인·편집 가능. \\
\bottomrule
\end{tabularx}
\end{center}

\begin{notebox}
  \textbf{DB 파일 백업}은 서버가 실행 중인 상태에서도 안전하게 수행됩니다.
  SQLite의 핫 스냅샷 API를 사용하므로 데이터 손상 없이 즉시 다운로드됩니다.
\end{notebox}

\subsubsection*{복원 방법}

\menu{관리자 대시보드 → 백업 / 복원} (\texttt{/admin/restore}) 페이지에서
두 가지 복원 방식을 모두 지원합니다.

\paragraph{Excel 파일 복원}
\begin{itemize}
  \item 복원할 항목(학생 계정·출결 기록·자습 신청·학습 기록)을 선택하여 업로드합니다.
  \item \textbf{추가 방식}으로 동작합니다 --- 기존 데이터를 덮어쓰거나 삭제하지 않습니다.
  \item 같은 파일을 두 번 업로드해도 중복이 생기지 않습니다.
  \item 신규로 생성되는 모든 계정에는 \textbf{랜덤 임시 비밀번호}가 일괄 발급되며,
        명단이 담긴 Excel(\texttt{self\_study\_복원결과\_...xlsx})이 복원 직후 자동 다운로드됩니다.
        학생·교사에게 전달 후 반드시 삭제하십시오.
  \item 자습실 QR 토큰은 새로 발급되므로 QR코드를 다시 인쇄·부착해야 합니다.
\end{itemize}

\paragraph{DB 파일 복원}
\begin{itemize}
  \item 이 시스템에서 내보낸 \texttt{.db} 파일을 업로드합니다.
  \item 현재 DB를 업로드한 파일로 \textbf{완전히 교체}합니다.
  \item 복원 전 기존 DB는 서버에 \texttt{self\_study.db.prev} 파일로 자동 보관됩니다.
  \item 업로드 파일은 메모리에 전부 올리지 않고 디스크로 스트리밍되며,
        매직 바이트와 스키마 검증 후에만 적용됩니다.
  \item 복원 직후 \textbf{관리자 본인 포함 모든 사용자의 세션이 무효화}됩니다.
        보안상 의도된 동작이며(\texttt{session\_token} 불일치로 자동 차단),
        새 DB의 비밀번호로 다시 로그인하면 됩니다.
  \item 옛 버전 DB(예: \texttt{system\_settings} 테이블 없음)로 복원해도
        직후 자동으로 누락 테이블이 생성되고 기본값이 시드되므로
        후속 페이지에서 \texttt{no such table} 같은 500 에러가 발생하지 않습니다.
\end{itemize}

\begin{notebox}
  업로드 크기 상한은 \textbf{50 MB}입니다 (서버 메모리 압박 방어).
  학교 단일 DB는 보통 수십 MB 이하라 충분하지만,
  더 큰 파일이 필요하면 \texttt{app.py}의 \texttt{MAX\_CONTENT\_LENGTH} 값을 조정하십시오.
  한도를 초과하면 브라우저에 ``413 Request Entity Too Large''가 표시됩니다.
\end{notebox}

\begin{warnbox}
  \textbf{DB 파일 복원}은 현재 운영 중인 모든 데이터를 덮어씁니다.
  복원 전에 반드시 현재 DB를 먼저 백업하십시오.
  업로드된 파일이 올바른 SQLite 파일인지 자동으로 검증하며,
  유효하지 않은 파일은 거부됩니다.
\end{warnbox}

\subsubsection*{시스템 업데이트 시 권장 절차}

\begin{enumerate}
  \item 관리자 대시보드에서 \textbf{DB 파일 백업(.db)} 다운로드
  \item (선택) \textbf{Excel 백업(.xlsx)}도 추가로 다운로드
  \item 코드 업데이트 수행 (git pull, 파일 교체 등)
  \item 필요한 경우 DB 마이그레이션 스크립트 실행 (아래 참조)
  \item 서버 재시작 후 정상 동작 확인
  \item 문제 발생 시 \menu{백업 / 복원} 페이지에서 백업한 \texttt{.db} 파일로 복원
\end{enumerate}

%========================================================
\subsection{DB 마이그레이션}
%========================================================

코드가 업데이트될 때 데이터베이스 구조(스키마)가 변경될 수 있습니다.
신규 설치에서는 서버를 처음 실행하면 최신 스키마로 자동 생성되지만,
\textbf{이미 운영 중인 DB}는 컬럼이나 제약 조건이 없는 상태일 수 있으므로
마이그레이션 스크립트를 별도로 실행해야 합니다.

\subsubsection*{마이그레이션 스크립트 목록}

\begin{center}
\begin{tabularx}{\textwidth}{llX}
\toprule
\textbf{파일} & \textbf{대상 버전} & \textbf{적용 내용} \\
\midrule
\texttt{migrate.py} & v1.9 이하 → 최신 &
  \texttt{attendance} 컬럼 추가(\texttt{checked\_out\_at}, \texttt{early\_leave\_note}),
  핵심 3개 테이블에 UNIQUE 제약 추가 \\
\texttt{migrate\_add\_settings.py} & 시스템 설정 도입 이전 → 최신 &
  \texttt{system\_settings} 테이블 생성 + 8개 운영 정책값 기본값 시드 \\
\texttt{migrate\_add\_constraints\_v2.py} & 데이터 무결성 제약 도입 이전 → 최신 &
  \texttt{users}·\texttt{attendance} 테이블 재구축으로 CHECK 제약(역할·성별·출결상태) 추가 \\
\texttt{migrate\_drop\_room\_seat\_uq.py} & v2 적용 후 → 최신 &
  \texttt{student\_rooms}의 \texttt{(study\_room\_id, seat\_number)} UNIQUE 제약 제거.
  남/여 zone 구조와 충돌해 정상 좌석 배정도 IntegrityError로 막던 문제 해결 \\
\bottomrule
\end{tabularx}
\end{center}

\begin{notebox}
  모든 스크립트는 \textbf{멱등(idempotent)}하게 작성되어 있습니다.
  이미 적용된 컬럼이나 제약이 있으면 건너뛰므로, 중복 실행해도 데이터가 손상되지 않습니다.
\end{notebox}

\subsubsection*{마이그레이션 실행 방법}

서버를 \textbf{중지한 상태}에서 서버 PC의 터미널(명령 프롬프트)을 열고 아래 명령을 실행합니다.

\begin{lstlisting}[language=bash, caption=마이그레이션 실행]
cd C:\Users\user\self_study    <- 운영 서버의 실제 경로로 변경
python migrate.py
\end{lstlisting}

정상 실행 시 다음과 유사한 출력이 나타납니다.

\begin{lstlisting}
[1/4] attendance 컬럼 확인...
  변경 없음.
[2/4] users 테이블 UNIQUE 제약 확인...
  이미 존재.
[3/4] attendance 테이블 UNIQUE 제약 확인...
  이미 존재.
[4/4] study_rooms 테이블 UNIQUE 제약 확인...
  이미 존재.

마이그레이션 완료.
\end{lstlisting}

처음 실행하는 경우 \texttt{이미 존재.} 대신 \texttt{컬럼 추가: ...} 또는 \texttt{UNIQUE 제약 추가 완료.} 메시지가 출력됩니다.

\subsubsection*{마이그레이션 적용 내용 상세}

\paragraph{(1) attendance 컬럼 추가}

기존 출결 테이블에 퇴실 관련 컬럼이 없으면 추가됩니다.

\begin{center}
\begin{tabularx}{\textwidth}{llX}
\toprule
\textbf{컬럼} & \textbf{타입} & \textbf{설명} \\
\midrule
\texttt{checked\_out\_at} & DATETIME & QR 퇴실 시각 \\
\texttt{early\_leave\_note} & VARCHAR(200) & 조퇴 사유 / 출석인정 메모 \\
\bottomrule
\end{tabularx}
\end{center}

\paragraph{(2) UNIQUE 제약 조건 추가}

중복 데이터 삽입을 DB 수준에서 차단하기 위해 다음 제약을 추가합니다.
기존 중복 데이터는 \texttt{INSERT OR IGNORE}로 처리되어 먼저 입력된 행이 유지됩니다.

\begin{center}
\begin{tabularx}{\textwidth}{llX}
\toprule
\textbf{테이블} & \textbf{제약 대상} & \textbf{의미} \\
\midrule
\texttt{users} & \texttt{student\_id} & 동일 학번 중복 가입 방지 \\
\texttt{attendance} & \texttt{(user\_id, date, period)} & 같은 날·교시 출결 중복 방지 \\
\texttt{study\_rooms} & \texttt{name} & 동일 이름 자습실 중복 방지 \\
\bottomrule
\end{tabularx}
\end{center}

\paragraph{(3) system\_settings 테이블 추가 (운영 정책값 도입)}

기존에는 사전 입실 시간(30분)·퇴실 grace(10분)·지각 임계(10분)·신청 마감일(20일)·
비밀번호 정책(8자, 영문+숫자) 등 학사 운영에 영향을 주는 값들이
소스 코드에 직접 박혀 있어 변경 때마다 코드 수정 + 서버 재시작이 필요했습니다.

\texttt{migrate\_add\_settings.py}는 다음 작업을 수행합니다.

\begin{enumerate}[leftmargin=1.5em]
    \item \texttt{system\_settings} 테이블을 생성합니다(이미 있으면 건너뜀).
    \item 8개 운영 정책값을 기본값(현재 운영 중인 값과 동일)으로 시드합니다.
\end{enumerate}

\begin{lstlisting}[language=bash, caption=시스템 설정 마이그레이션 실행]
cd C:\Users\user\self_study    <- 운영 서버의 실제 경로로 변경
python migrate_add_settings.py
\end{lstlisting}

정상 실행 시 다음과 유사한 출력이 나타납니다.

\begin{lstlisting}
system_settings 테이블 생성 완료.
기본 설정값 8개 시드 완료.
\end{lstlisting}

\begin{notebox}
  이 스크립트는 \textbf{기본값을 현재 운영 중이던 값과 동일하게} 시드하므로,
  실행 직후 학생·교사가 체감하는 동작은 전혀 바뀌지 않습니다.
  실제 정책 변경은 \menu{관리자 대시보드 → 시스템 설정}에서 수행합니다.
\end{notebox}

\begin{notebox}
  앱이 한 번이라도 부팅되었다면 시드는 자동으로 수행됩니다.
  이 스크립트는 \textbf{서버를 멈춘 상태에서 미리 적용}하고 싶거나
  자동 시드가 동작하지 않는 경우의 안전장치입니다.
\end{notebox}

\paragraph{(4) 데이터 무결성 제약 추가 (CHECK · UNIQUE)}

\texttt{migrate\_add\_constraints\_v2.py}는 다음 제약을 DB에 추가합니다.
SQLite는 \texttt{ALTER TABLE}로 제약을 직접 추가할 수 없으므로 각 테이블을
임시 테이블로 재구축한 뒤 이름을 바꿔주는 방식으로 동작합니다.

\begin{center}
\begin{tabularx}{\textwidth}{llX}
\toprule
\textbf{테이블} & \textbf{제약} & \textbf{의미} \\
\midrule
\texttt{users}
  & CHECK \texttt{role}
  & \texttt{student} / \texttt{teacher} / \texttt{admin} 이외 값 차단 \\
\texttt{users}
  & CHECK \texttt{gender}
  & NULL 또는 \texttt{M} / \texttt{F} 이외 값 차단 \\
\texttt{attendance}
  & CHECK \texttt{status}
  & \texttt{present, late, absent, early\_leave, approved\_leave, after\_school} 이외 차단 \\
\bottomrule
\end{tabularx}
\end{center}

\begin{notebox}
  초기 v2에는 \texttt{student\_rooms (study\_room\_id, seat\_number)} UNIQUE 제약도 포함됐으나,
  자습실의 남/여 zone 구조에서 ``남학생 1번''과 ``여학생 1번''이 별도의 물리 좌석을 의미하는
  현실과 충돌해 \texttt{migrate\_drop\_room\_seat\_uq.py}로 제거되었습니다.
  zone 내 좌석 번호 중복은 코드 레벨(\texttt{random.sample}, 수동 배정의 \texttt{used\_m}/\texttt{used\_f}
  set)이 담당합니다.
\end{notebox}

\subparagraph*{실행 순서}

\begin{enumerate}[leftmargin=1.5em]
  \item 스크립트가 먼저 \textbf{사전 검사}를 돌려 현재 데이터가 위 제약을 이미 위반하는지 스캔합니다.
        위반이 있으면 상세 목록을 출력하고 중단합니다 --- 강제 실행 옵션은 없습니다.
  \item 위반이 없으면 트랜잭션 내에서 3개 테이블을 순서대로 재구축합니다.
  \item 완료 후 \texttt{PRAGMA foreign\_key\_check}로 FK 무결성을 재확인합니다.
\end{enumerate}

\begin{lstlisting}[language=bash, caption=제약 마이그레이션 실행]
cd C:\Users\user\self_study    <- 운영 서버의 실제 경로로 변경
python migrate_add_constraints_v2.py
\end{lstlisting}

정상 실행 시 출력 예:

\begin{lstlisting}
[1/3] 사전 검사 통과 - 제약 위반 데이터 없음.

[2/3] 테이블 재구축 (트랜잭션)...
  - users: 재구축 시작...
  - users: 완료
  - attendance: 재구축 시작...
  - attendance: 완료
  - student_rooms: 재구축 시작...
  - student_rooms: 완료

[3/3] 외래 키 무결성 검사...
  OK

마이그레이션 완료. 3개 테이블 재구축.
\end{lstlisting}

\subparagraph*{사전 검사 실패 시}

\texttt{attendance}에 옛 알 수 없는 \texttt{status} 값, \texttt{student\_rooms}에 좌석 번호 중복 등이 발견되면
스크립트가 중단되고 해당 행을 출력합니다.
관리자 UI 또는 직접 SQL로 정정한 뒤 재실행하십시오. 중복 좌석은 한 쪽의 \texttt{seat\_number}를
\texttt{NULL}로 변경하면 제약을 피하면서 배정 관계는 유지할 수 있습니다.

\subparagraph*{부팅 시 자동 감지 (안전망)}

\texttt{models.py}에 \texttt{CheckConstraint}/\texttt{UniqueConstraint}를 추가해도
SQLAlchemy의 \texttt{db.create\_all()}은 \textbf{기존 테이블을 변경하지 않습니다}.
즉, 코드만 업데이트하고 마이그레이션을 실행하지 않으면 운영 DB에는 제약이 없는 상태가 유지됩니다.

이 위험을 가시화하기 위해 앱 부팅 시 4개 핵심 제약(\texttt{ck\_users\_role},
\texttt{ck\_users\_gender}, \texttt{ck\_attendance\_status}, \texttt{uq\_room\_seat})의 적용 여부를
자동 검사합니다. 누락 시:

\begin{itemize}[leftmargin=1.5em]
  \item \texttt{logs/audit.log}에 \texttt{system.constraints\_missing} 이벤트 기록 (WARNING)
  \item 콘솔에 \texttt{!!!} 경고 박스 출력 (\texttt{logs/app.log}에도 남음)
  \item 앱은 \textbf{정상 부팅} (운영 중단 X) --- 관리자가 마이그레이션을 수행하도록 유도
\end{itemize}

따라서 코드 업데이트 후 운영 PC를 재시작했을 때 \texttt{logs/audit.log}나 콘솔에
이 경고가 보이면 \textbf{즉시} \texttt{python migrate\_add\_constraints\_v2.py}를 실행하십시오.

\begin{warnbox}
  이 마이그레이션은 \textbf{테이블을 완전히 재구축}합니다.
  사전 검사를 통과해도 드물게 FK 이름 꼬임이 발생할 수 있으므로,
  실행 전 반드시 \textbf{DB 파일 백업}을 받은 뒤 진행하십시오.
\end{warnbox}

\begin{warnbox}
  마이그레이션 실행 전에 반드시 \textbf{DB 파일을 백업}하십시오.
  오류 발생 시 자동 롤백되지만, 백업이 있으면 더 안전하게 복구할 수 있습니다.
\end{warnbox}

\subsubsection*{마이그레이션 실패 시 복구}

스크립트 실행 중 오류가 발생하면 트랜잭션이 자동으로 롤백되고 오류 메시지가 출력됩니다.

\begin{lstlisting}
오류 발생, 롤백됨: <오류 내용>
\end{lstlisting}

이 경우 다음 순서로 복구합니다.

\begin{enumerate}
  \item 출력된 오류 메시지를 확인합니다.
  \item 관리자 화면의 \menu{백업 / 복원}에서 사전에 백업한 \texttt{.db} 파일을 복원합니다.
  \item 오류 원인을 해결한 후 스크립트를 다시 실행합니다.
\end{enumerate}

%========================================================
\section{자주 묻는 질문 (FAQ)}
%========================================================

\subsection*{Q1. 교사로 회원가입했는데 로그인이 안 됩니다.}

교사 계정은 관리자의 승인 후에만 로그인할 수 있습니다.
관리자에게 계정 승인을 요청하십시오.

\subsection*{Q2. 학생이 자습을 신청하지 않으면 QR 출석이 안 됩니까?}

예. 해당 날짜·교시에 자습 신청을 완료한 학생만 QR 출석이 가능합니다.
신청하지 않은 학생이 QR코드를 스캔하면 ``신청하지 않았습니다''라는 안내가 표시됩니다.
교사는 출석 관리 페이지에서 직접 출석을 기록할 수 있습니다.

\subsection*{Q2-1. 학생이 퇴실 QR을 스캔했는데 ``입실 기록이 없습니다''라고 합니다.}

퇴실 QR은 반드시 같은 날 같은 교시에 입실 QR을 먼저 스캔한 후에만 동작합니다.
입실 QR을 스캔하지 않은 상태에서 퇴실 QR을 스캔하면 이 메시지가 표시됩니다.

\subsection*{Q2-2. 조퇴 처리된 학생을 출석인정으로 바꾸려면 어떻게 합니까?}

\menu{출석 관리} 페이지에서 조퇴 학생의 셀에 \menu{사유 입력} 필드와
\menu{출석인정} 버튼이 자동으로 표시됩니다.
사유(선택)를 입력하고 \menu{출석인정} 버튼을 클릭하면 상태가 변경되고
변경 이력에 기록됩니다.

\subsection*{Q3. QR코드 스캔 후 아무 반응이 없습니다.}

스마트폰 카메라 앱의 QR코드 기능을 사용하거나, QR코드 전용 앱을 이용하십시오.
URL이 자동으로 열리지 않으면 브라우저에 직접 URL을 복사하여 접속하십시오.

\subsection*{Q4. 학생이 다른 달의 자습을 신청할 수 있습니까?}

이전·다음 달로 이동하여 신청할 수 있습니다.
단, 20일 이전이면 기본적으로 다음 달이 표시됩니다.

\subsection*{Q5. 공휴일에 자습이 없도록 설정하려면?}

공휴일을 등록해도 자습은 가능합니다.
공휴일 자습을 없애려면 \menu{설정 → 공휴일} 탭에서 모든 교시의 \textbf{활성화 체크박스를 해제}하십시오.

\subsection*{Q6. 학생·교사가 비밀번호를 잊어버렸습니다.}

관리자가 \menu{사용자 관리} 페이지에서 해당 계정의 \menu{PW 초기화} 버튼을 클릭하면
임시 비밀번호가 생성됩니다. 임시 비밀번호를 학생·교사에게 전달하고
로그인 후 즉시 변경하도록 안내하십시오.

\subsection*{Q7. 관리자 비밀번호를 잊어버렸습니다.}

5.6절의 ``비밀번호 분실 시 재설정'' 방법을 참고하십시오.

\subsection*{Q7-1. \texttt{logs/app.log}와 \texttt{logs/audit.log}는 무엇이 다릅니까?}

\begin{itemize}[leftmargin=1.5em]
  \item \texttt{logs\textbackslash app.log}: 서버의 표준 출력입니다.
        Waitress가 받은 모든 HTTP 요청, 파이썬 예외, 시작·종료 메시지가 들어갑니다.
        주로 \textbf{기술적 장애 진단}에 씁니다.
  \item \texttt{logs\textbackslash audit.log}: 5.4절의 11개 추적 이벤트만
        구조화된 형태로 기록됩니다. 5MB 단위로 회전(\texttt{audit.log.1}\textasciitilde\texttt{audit.log.5})하며,
        ``누가 언제 어떤 관리자 작업을 했는가''를 \textbf{운영 책임 추적}에 씁니다.
\end{itemize}

평문 비밀번호는 어느 쪽에도 기록되지 않습니다.

\subsection*{Q7-2. DB 백업 파일 업로드 시 ``413 Request Entity Too Large''가 표시됩니다.}

업로드 크기 상한은 50 MB입니다. 학교 단일 DB는 보통 그보다 훨씬 작지만,
오랜 운영 후 출결·학습 기록이 누적되어 한도를 넘었다면
\texttt{app.py}의 \texttt{MAX\_CONTENT\_LENGTH} 값(바이트 단위)을 조정한 뒤 서버를 재시작하십시오.

\subsection*{Q8. 여러 대의 컴퓨터에서 동시에 사용할 수 있습니까?}

가능합니다. 서버는 \texttt{host='0.0.0.0'}으로 실행되어 학내 Wi-Fi(교사망·학생망 모두)에
연결된 모든 기기에서 서버 PC의 IP 주소로 동시 접속할 수 있습니다.

\subsection*{Q9. Wi-Fi가 끊겼다가 재연결되면 서버가 응답하지 않습니다.}

Flask 내장 개발 서버는 네트워크 변동 시 소켓이 깨지는 문제가 있습니다.
본 시스템은 \textbf{Waitress} WSGI 서버를 사용하도록 설정되어 있어 Wi-Fi 재연결 후에도
자동으로 복구됩니다. 3.5절의 자동 시작 설정을 완료하면 재부팅 후에도 자동으로 실행됩니다.
만약 문제가 계속된다면 작업 스케줄러의 ``실패 시 자동 재시작'' 옵션이 활성화되어 있는지 확인하십시오.

\subsection*{Q10. 서버를 재시작하면 로그인이 풀립니까?}

아닙니다. 세션 암호화 키가 \texttt{instance/secret\_key.txt}에 저장되므로
서버를 재시작해도 기존 로그인 세션이 유지됩니다.
단, \texttt{secret\_key.txt}를 삭제하면 새 키가 생성되어 모든 사용자가 로그아웃됩니다.

%========================================================
\section{URL 목록}
%========================================================

\begin{center}
\begin{tabularx}{\textwidth}{ll>{\raggedright\arraybackslash}X}
\toprule
\textbf{역할} & \textbf{URL} & \textbf{기능} \\
\midrule
공통 & \texttt{/login} & 로그인 \\
공통 & \texttt{/register} & 회원가입 \\
공통 & \texttt{/logout} & 로그아웃 \\
\midrule
관리자 & \texttt{/admin/} & 관리자 대시보드 \\
관리자 & \texttt{/admin/users} & 사용자 관리 (학생·교사 목록) \\
관리자 & \texttt{/admin/users/<ID>/reset-password} & 비밀번호 초기화 (POST) \\
관리자 & \texttt{/admin/users/<ID>/delete} & 계정 삭제 (POST) \\
관리자 & \texttt{/admin/teachers} & 교사 승인 관리 \\
관리자 & \texttt{/admin/change-password} & 마이페이지 (비밀번호 변경) \\
관리자 & \texttt{/admin/new-year} & 새 학년도 초기화 \\
관리자 & \texttt{/admin/new-year/backup} & 학년도 종료 전 Excel 백업 다운로드 \\
관리자 & \texttt{/admin/db-backup} & DB 파일 전체 백업 다운로드 (.db) \\
관리자 & \texttt{/admin/restore} & 백업 Excel 파일로 DB 복원 \\
관리자 & \texttt{/admin/db-restore} & DB 파일로 완전 복원 (POST) \\
\midrule
교사 & \texttt{/teacher/} & 교사 대시보드 \\
교사 & \texttt{/teacher/students} & 학생 목록·공간 배정 \\
교사 & \texttt{/teacher/applications} & 자습 신청 현황 \\
교사 & \texttt{/teacher/attendance} & 출석 관리 \\
교사 & \texttt{/teacher/attendance/auto-process} & 자동 지각/결석 처리 (POST) \\
교사 & \texttt{/teacher/attendance/update} & 출석 상태 수정 (POST) \\
교사 & \texttt{/teacher/attendance/export} & 출석 현황 Excel 다운로드 \\
교사 & \texttt{/teacher/statistics} & 참여 통계 (총 자습 시간 포함) \\
교사 & \texttt{/teacher/statistics/export} & 참여 통계 Excel 다운로드 (총 자습 시간 포함) \\
교사 & \texttt{/teacher/mypage} & 마이페이지 (내 정보·비밀번호 변경) \\
교사 & \texttt{/teacher/student/<ID>/report} & 학생 개인 리포트 \\
교사 & \texttt{/teacher/settings} & 자습 시간·공간·공휴일 설정 \\
교사 & \texttt{/teacher/qr/<공간ID>} & QR코드 표시 \\
교사 & \texttt{/teacher/seat-layout/<공간ID>} & 좌석 배치도 (드래그 앤 드롭 + 실시간 출결) \\
교사 & \texttt{/teacher/seat-layout/<공간ID>/save} & 배치도 저장 (POST, JSON) \\
교사 & \texttt{/teacher/api/room/<공간ID>/attendance\_status} & 자습실 출결 현황 JSON API \\
\midrule
학생 & \texttt{/student/} & 학생 대시보드 \\
학생 & \texttt{/student/apply} & 월별 자습 신청 \\
학생 & \texttt{/student/my-attendance} & 내 출결 현황 (월별 자기확인) \\
학생 & \texttt{/student/log} & 학습 기록 \\
학생 & \texttt{/student/qr-attend/<토큰>} & QR 입실 (출석) 처리 \\
학생 & \texttt{/student/qr-checkout/<토큰>} & QR 퇴실 (조퇴) 처리 \\
학생 & \texttt{/student/mypage} & 마이페이지 (내 정보·비밀번호 변경) \\
\bottomrule
\end{tabularx}
\end{center}

\vfill
\begin{center}
\rule{0.6\textwidth}{0.4pt}\\[0.5em]
{\small 본 문서는 자율학습 관리 시스템 v1.9 기준으로 작성되었습니다.}\\[0.2em]
{\small v1.2 변경: QR 퇴실·조퇴 처리·출석인정 기능 추가}\\[0.1em]
{\small v1.3 변경: 관리자 DB 파일 백업(.db) 및 복원 기능 추가}\\[0.1em]
{\small v1.4 변경: 자습 공간 수정 기능 추가, 교사 계정·자습실 배정 백업/복원 추가}\\[0.1em]
{\small v1.5 변경: 자습실 목록(QR 토큰 포함) 백업/복원 추가, 학생 비밀번호 백업/복원 수정}\\[0.1em]
{\small v1.6 변경: 공휴일·자습시간설정·교사 담당학년 백업/복원 추가}\\[0.1em]
{\small v1.7 변경: 관리자 계정 백업/복원 추가, 백업 시트 순서 정리 (11개 시트 완전 백업)}\\[0.1em]
{\small v1.8 변경: 평일 아침 자습(07:30--08:30) 신설, 교사·학생 마이페이지 추가, 참여 통계에 총 자습 시간 표시}\\[0.1em]
{\small v1.9 변경: 요일별 자습 시간 개별 설정, QR 사전 입실(30분 전), 퇴실 유예(10분) 및 재스캔 시 최종 시간 반영}\\[0.1em]
{\small v2.0 변경: 통합 마이그레이션 스크립트(\texttt{migrate.py}) 추가, SQLite FK 강제·WAL 모드 활성화, QR 입실 확인 화면 추가}
\end{center}

\end{document}
